% Pulsar Technical Specification --- NIST MPTC submission
% Target deadline: 2026-Nov-16
%
% Compile: latexmk -pdf pulsar.tex
%
% Status: cryptographic body drafted. Encoding section is intentionally
% structural only --- byte-level wire formats freeze at DD-008
% (end of August 2026).

\documentclass[11pt,a4paper]{article}

\usepackage[margin=1in]{geometry}
\usepackage{amsmath, amssymb, amsthm}
\usepackage{algorithm, algpseudocode}
\usepackage{hyperref}
\usepackage{cite}
\usepackage{microtype}
\usepackage{xcolor}
\usepackage{listings}
\usepackage{booktabs}
\usepackage{enumitem}
\usepackage{cleveref}

\hypersetup{
  colorlinks=true,
  linkcolor=blue!50!black,
  citecolor=blue!50!black,
  urlcolor=blue!50!black,
}

\newtheorem{theorem}{Theorem}
\newtheorem{lemma}[theorem]{Lemma}
\newtheorem{proposition}[theorem]{Proposition}
\newtheorem{corollary}[theorem]{Corollary}
\theoremstyle{definition}
\newtheorem{definition}{Definition}
\newtheorem{game}{Game}
\newtheorem{construction}{Construction}
\theoremstyle{remark}
\newtheorem{remark}{Remark}
\newtheorem{sketch}{Proof sketch}
\newtheorem{openq}{Open question}

\crefname{theorem}{Theorem}{Theorems}
\crefname{lemma}{Lemma}{Lemmas}
\crefname{proposition}{Proposition}{Propositions}
\crefname{corollary}{Corollary}{Corollaries}
\crefname{definition}{Definition}{Definitions}
\crefname{game}{Game}{Games}
\crefname{construction}{Construction}{Constructions}
\crefname{remark}{Remark}{Remarks}
\crefname{openq}{Open question}{Open questions}
\crefname{algorithm}{Algorithm}{Algorithms}

\newcommand{\Rq}{R_q}
\newcommand{\Zq}{\mathbb{Z}_q}
\newcommand{\Tq}{\mathbb{T}_q}
\newcommand{\GF}{\mathrm{GF}}
\newcommand{\negl}{\mathsf{negl}}
\newcommand{\poly}{\mathsf{poly}}
\newcommand{\PPT}{\mathsf{PPT}}
\newcommand{\Adv}{\mathsf{Adv}}
% \Pr is provided by LaTeX as a math operator; we use it directly.
\newcommand{\MLWE}{\mathsf{MLWE}}
\newcommand{\MSIS}{\mathsf{MSIS}}
\newcommand{\SHAKE}{\mathrm{SHAKE256}}
\newcommand{\cSHAKE}{\mathrm{cSHAKE256}}
\newcommand{\KMAC}{\mathrm{KMAC256}}
\newcommand{\TupleHash}{\mathrm{TupleHash256}}
\newcommand{\NTT}{\mathrm{NTT}}
\newcommand{\INTT}{\mathrm{NTT}^{-1}}
\newcommand{\HighBits}{\mathsf{HighBits}}
\newcommand{\LowBits}{\mathsf{LowBits}}
\newcommand{\MakeHint}{\mathsf{MakeHint}}
\newcommand{\UseHint}{\mathsf{UseHint}}
\newcommand{\Power}{\mathsf{Power2Round}}
\newcommand{\Decompose}{\mathsf{Decompose}}
\newcommand{\Sign}{\textsc{Sign}}
\newcommand{\Verify}{\textsc{Verify}}
\newcommand{\Setup}{\textsc{Setup}}
\newcommand{\DKG}{\textsc{DKG}}
\newcommand{\Reshare}{\textsc{Reshare}}
\newcommand{\Refresh}{\textsc{Refresh}}
\newcommand{\Comm}{\mathsf{Comm}}

\title{Pulsar:\\
       \large A Threshold Implementation of FIPS 204 ML-DSA\\
       NIST MPTC Class N1 + N4 Candidate}
\author{Lux Industries Inc.\\
        \texttt{security@lux.network}}
\date{Draft \today}

\begin{document}
\maketitle

\begin{abstract}
Pulsar is a $t$-of-$n$ threshold implementation of NIST FIPS 204
(ML-DSA)~\cite{nist-fips204} signing and distributed key generation.
The threshold ceremony emits signatures byte-equal to single-party
FIPS~204 ML-DSA signatures: Pulsar signatures verify under
unmodified \texttt{ML-DSA.Verify}. The protocol structure is the
two-round threshold scheme of Boschini, Kaviani, Lai, Malavolta,
Takahashi, and Tibouchi (``Corona''~\cite{boschini2024corona}), retargeted
from Ring-LWE ($\Rq$, single polynomial) to Module-LWE ($\Rq^k$,
polynomial vectors) so the protocol's algebraic shape matches
FIPS~204.

This document is the technical specification submitted to the NIST
Multi-Party Threshold Cryptography first call (NIST IR
8214C,~\cite{nist-ir-8214c}). We claim Class~N1 (signing) and Class~N4
(ML keygen / DKG): output interchangeability with FIPS~204 is the
headline contribution.

\paragraph{Status.} Draft. Encoding freeze targeted for end-August
2026; package submission targeted for 2026-Nov-16. Sections on
encodings and concrete benchmarks are intentionally structural only
in this draft --- the byte-level wire formats freeze at the
encoding-freeze gate.
\end{abstract}

\tableofcontents

% =====================================================================
\section{Introduction}
% =====================================================================

\paragraph{FIPS~204 ML-DSA.}
FIPS~204~\cite{nist-fips204} standardises the Module-Lattice-Based
Digital Signature Algorithm (ML-DSA), a derandomised variant of
CRYSTALS-Dilithium. ML-DSA is a Fiat--Shamir-with-Aborts scheme over
the cyclotomic ring $\Rq = \Zq[X]/(X^{256}+1)$ with prime modulus
$q = 2^{23} - 2^{13} + 1 = 8\,380\,417$. Three parameter sets ---
ML-DSA-44, ML-DSA-65, ML-DSA-87 --- target NIST PQ categories
$2$, $3$, and $5$ respectively. The signature format is a triple
$\sigma = (\tilde c, z, h)$ where $\tilde c \in \{0,1\}^{2\lambda}$
is the commitment hash, $z \in \Rq^\ell$ is the masked response, and
$h \in \{0,1\}^{\omega + k}$ encodes hint bits.

\paragraph{Threshold signing.}
A $t$-of-$n$ threshold signature scheme replaces the single signing
party with a quorum: any subset of at least $t$ parties from a fixed
group of $n$ can jointly produce a signature, while any subset of
size strictly less than $t$ learns no useful information about the
secret key and cannot forge. Threshold deployments are mandatory for
high-value systems where single-key custody is unacceptable
(consensus validators, custodial-key cold storage, root-CA signing).
NIST IR~8214C~\cite{nist-ir-8214c} structures the design space into
\emph{Class~S} (special threshold-friendly primitives, no output
interchangeability requirement) and \emph{Class~N} (threshold
implementations \emph{of} a specified NIST primitive whose outputs
are interchangeable with the non-threshold primitive's outputs).

\paragraph{Output interchangeability is the headline claim.}
The Class~N positioning is strictly stronger: a threshold signature
produced by Pulsar verifies against the same group public key
$\mathit{pk}$ under unmodified \texttt{ML-DSA.Verify}. Any deployed
FIPS-validated ML-DSA verifier --- BoringSSL FIPS, AWS-LC, OpenSSL~3,
liboqs, or any future module conforming to FIPS~204 --- consumes
Pulsar signatures with no code change. This collapses the
threshold/non-threshold split at the wire format: a relying party
who can verify ML-DSA can verify Pulsar.

\paragraph{Why Pulsar alongside Pulsar.}
The companion scheme Corona (R-LWE)~\cite{boschini2024corona} is a
Class~S candidate over the single-polynomial ring $\Rq$. Pulsar's
construction is faster, simpler, and produces a different wire
format (the verifier is \emph{not} FIPS~204; it is a Pulsar-specific
verifier). Pulsar ports that protocol's $2$-round signing structure
to the Module-LWE setting $\Rq^k$ in which FIPS~204 operates so that
the output byte string is exactly an ML-DSA signature. We retain
Pulsar's Pedersen DKG (over $\Rq^k$ rather than $\Rq$) and Pulsar's
HJKY97-style proactive resharing, modulo:

\begin{itemize}[leftmargin=*]
\item \textbf{DD-002 (hash family).} The NIST profile uses
exclusively FIPS~202~\cite{nist-fips202} +
SP~800-185~\cite{nist-sp800-185}: $\SHAKE$, $\cSHAKE$, $\KMAC$,
$\TupleHash$. No BLAKE3. The same SHAKE-256 sponge is used
internally by ML-DSA-65 (FIPS~204~\S5.1).
\item \textbf{DD-003 (lattice shape).} The protocol operates on
$\Rq^k$ with the FIPS~204 parameters $(k, \ell, \eta, \tau, \gamma_1,
\gamma_2, \omega, \beta, d, \lambda)$ inherited verbatim per
parameter level.
\item \textbf{DD-006 (reshare quorum selection).} Long-lived
deployments select the reshare quorum from the chain's randomness
beacon rather than deterministically (smallest-ID). The deterministic
selection used by the Pulsar reference implementation lets a mobile
adversary game the corrupt-set positioning across epochs; the beacon
selection blocks that line of attack.
\end{itemize}

\paragraph{High-level construction.}
Pulsar consists of four protocols and one verifier:

\begin{itemize}[leftmargin=*]
\item $\DKG$: Pedersen distributed key generation over $\Rq^k$.
Three rounds. Output: a group public key $\mathit{pk}$ distributed
exactly as a fresh ML-DSA public key, and a per-party Shamir share
$(s_i, u_i)$ where $s_i$ shares the secret $s$ under the FIPS~204
secret distribution and $u_i$ is the blinding share used for
Pedersen binding.
\item $\Sign$: Two-round threshold signing. Inputs: shares
$(s_i, u_i)$ and message $\mu$. Output: a single FIPS~204
$\sigma = (\tilde c, z, h)$ that verifies under \texttt{ML-DSA.Verify}.
\item $\Verify$: literally \texttt{ML-DSA.Verify} from FIPS~204
\S6.3; we quote and \emph{do not modify} the verifier.
\item $\Reshare$: HJKY97-style proactive resharing, with
beacon-randomised quorum selection. Inputs: old committee shares,
new committee identity set, fresh beacon randomness.
Output: new committee shares producing the \emph{same} $\mathit{pk}$.
\item Identifiable abort: every protocol stage emits signed evidence
of any detected deviation, suitable for slashing.
\end{itemize}

\paragraph{Document map.}
\S\ref{sec:notation} pins notation. \S\ref{sec:prelim} states the
hardness assumption (M-LWE), Pedersen commitments over $\Rq^k$,
Shamir secret sharing over $\Rq^k$, and the FIPS~204 sampling
primitives we re-use verbatim. \S\ref{sec:algorithms} specifies
$\DKG$, $\Sign$, the verifier, and $\Reshare$.
\S\ref{sec:security} defines the security games and reduces each to
either M-LWE / MSIS or to the FIPS~204 EUF-CMA argument; the formal
game definitions live in the companion file
\texttt{spec/security-games.tex} included verbatim.
\S\ref{sec:parameters} specifies three parameter sets matching
ML-DSA-44/65/87 with lattice-estimator commitments.
\S\ref{sec:encodings} is structural only in this draft; byte-level
wire formats freeze at DD-008.
\S\ref{sec:system-model}--\ref{sec:deployment} cover the system model,
implementation considerations, and operational deployment guidance.

% =====================================================================
\section{Notation}
\label{sec:notation}
% =====================================================================

\paragraph{Rings.} For a prime $q$ and positive power-of-two $n$, let
$\Rq = \Zq[X]/(X^n + 1)$ be the cyclotomic ring of degree $n$ over
$\Zq$. Pulsar operates over $\Rq$ with $(q, n) = (8\,380\,417, 256)$
matching FIPS~204. Lower-case Latin letters denote ring elements;
bold or vector arrows denote elements of $\Rq^k$ or $\Rq^\ell$.

\paragraph{NTT.} The number-theoretic transform
$\NTT : \Rq \to \Zq^n$ is the FIPS~204 NTT
(FIPS~204~\S7.5). $\NTT$ is a ring isomorphism; for $a, b \in \Rq$,
$\NTT(a \cdot b) = \NTT(a) \odot \NTT(b)$ where $\odot$ denotes
coordinate-wise multiplication in $\Zq^n$. Multiplication of ring
elements is computed by transporting both factors to NTT form,
multiplying coordinate-wise, and inverse-transforming. Throughout
this specification we adopt the convention that all matrix and
vector multiplications are computed in NTT representation unless
explicitly noted; this is the same convention as FIPS~204.

\paragraph{Module shape.} The FIPS~204 ML-DSA parameter sets are
indexed by a pair $(k, \ell)$ giving the dimensions of the public
matrix $\mathbf{A} \in \Rq^{k \times \ell}$ and the secret vector
$\mathbf{s}_1 \in \Rq^\ell$. The three FIPS~204 parameter sets fix
$(k, \ell) \in \{(4, 4), (6, 5), (8, 7)\}$ for ML-DSA-44, ML-DSA-65,
ML-DSA-87 respectively (FIPS~204~\S4 Table~1).

\paragraph{Threshold parameters.} A threshold scheme has total party
count $n$ and reconstruction threshold $t$, with $1 \le t \le n$ and
$n \ge 2t - 1$ (so a $t$-element quorum exists alongside up to $t-1$
faulty parties). Each party is identified by a $1$-indexed integer
$i \in \{1, \dots, n\}$. We forbid the index $0$ because $0$ is the
Shamir evaluation point of the secret. We re-use the symbol $n$ for
both the ring degree and the party count where context disambiguates;
where a clash would be confusing we write $\nu$ for the ring degree
(matching FIPS~204's variable name).

\paragraph{Hash family.} All hashing uses the FIPS~202 +
SP~800-185 family: $\SHAKE$ (FIPS~202~\S6.3), $\cSHAKE$
(SP~800-185~\S3), $\KMAC$ (SP~800-185~\S4), $\TupleHash$
(SP~800-185~\S5). Each call is parameterised by a customisation
string. The canonical customisation tags used by Pulsar are:

\begin{center}\small
\begin{tabular}{l l}
\toprule
purpose & SP~800-185 customisation tag \\
\midrule
DKG commit-digest (Round~1.5)            & \texttt{PULSAR-DKG-COMMIT-V1} \\
DKG transcript hash                       & \texttt{PULSAR-DKG-TRANSCRIPT-V1} \\
Sign Round-1 transcript (sid hash $H_1$)  & \texttt{PULSAR-SIGN-R1-V1} \\
Sign Round-2 challenge (FIPS~204 $\tilde c$) & FIPS~204 inner SHAKE-256 \\
Sign per-party PRNG (round-keyed)         & \texttt{PULSAR-SIGN-PRNG-V1} \\
Sign PRF for mask cancellation            & \texttt{PULSAR-SIGN-PRF-V1} \\
Reshare commit digest                      & \texttt{PULSAR-RESHARE-COMMIT-V1} \\
Reshare transcript                         & \texttt{PULSAR-RESHARE-TRANSCRIPT-V1} \\
Reshare beacon binding                     & \texttt{PULSAR-RESHARE-BEACON-V1} \\
Public matrix $\mathbf{A}$ derivation       & FIPS~204 \texttt{ExpandA} (no change) \\
Public matrix $\mathbf{B}$ derivation       & \texttt{PULSAR-EXPANDB-V1} \\
Complaint signing transcript               & \texttt{PULSAR-COMPLAINT-V1} \\
\bottomrule
\end{tabular}
\end{center}

The bound on customisation-string length is governed by
SP~800-185~\S3.2 ($\le 2^{2040}$ bits); all Pulsar tags above are
short-form ASCII strings well within the bound. Version
suffixes (\texttt{-V1}) are reserved for future rotation; rotating any
tag invalidates every test vector pinned at that tag.

\paragraph{Sampling.} Five distributions appear in this specification:

\begin{itemize}[leftmargin=*]
\item \emph{Uniform mod $q$}, written $a \gets \Zq$ (per coefficient
slot). The FIPS~204 rejection-based uniform sampler is reused
verbatim; see FIPS~204~\S7.3.
\item \emph{Centred binomial $B_\eta$} over the integers, written
$a \gets B_\eta$. Secret coefficients have this distribution;
FIPS~204~\S7.4 (\texttt{SampleInBall} is unrelated --- the secret
sampler is the centred binomial of \S7.4 with the parameter
$\eta$ given per parameter set).
\item \emph{Ternary with Hamming weight $\tau$}, written
$c \gets B_\tau$. Used for the challenge polynomial $c$; FIPS~204
\texttt{SampleInBall} (\S7.6).
\item \emph{Discrete Gaussian with parameter $\sigma$}, written
$a \gets D_\sigma$. Used only inside the Pedersen DKG round-1
sampling; see \S\ref{sec:prelim-gauss}.
\item \emph{Uniform interval} $[-\gamma_1 + 1, \gamma_1]$, written
$y \gets U_{\gamma_1}$. Used for masking the response polynomial in
the FIPS~204 signer; FIPS~204~\S7.7.
\end{itemize}

\paragraph{Norms.} Coefficients of $\Zq$ are canonically lifted to
$(-q/2, q/2]$ before any infinity-norm or two-norm is taken; we
write $\|a\|_\infty$ and $\|a\|_2$ throughout. For a vector
$\mathbf{v} \in \Rq^k$, $\|\mathbf{v}\|_\infty = \max_i \|v_i\|_\infty$,
the maximum entry over all coefficient slots of all components.

\paragraph{FIPS~204 macros.} We retain the FIPS~204 names
$\HighBits, \LowBits, \Decompose, \Power, \MakeHint, \UseHint$ for the
high-bits and hint extraction routines (FIPS~204~\S7.4--7.5); their
behaviour and parameterisation by $\gamma_2$ and $d$ are inherited
unchanged.

% =====================================================================
\section{Mathematical preliminaries}
\label{sec:prelim}
% =====================================================================

\subsection{Module-LWE assumption}
\label{sec:prelim-mlwe}

The hardness assumption underlying Pulsar is identical to the
hardness assumption underlying FIPS~204 ML-DSA: the decisional
Module-LWE problem over $\Rq^k$.

\begin{definition}[Decisional Module-LWE]
\label{def:mlwe}
Fix $(q, n, k, \ell, \eta) \in \mathbb{N}^5$. Let $\chi_\eta$ denote the
centred binomial distribution of parameter $\eta$ over $\Zq$, applied
coefficient-wise to elements of $\Rq$. The decisional
$\MLWE_{q, k, \ell, \eta}$ problem is to distinguish, with
non-negligible advantage, the distribution
\[
\bigl(\mathbf{A},\,\mathbf{A} \cdot \mathbf{s}_1 + \mathbf{s}_2\bigr)
\quad \text{where } \mathbf{A} \gets \Rq^{k \times \ell},\;
\mathbf{s}_1 \gets \chi_\eta^\ell,\;
\mathbf{s}_2 \gets \chi_\eta^k,
\]
from the uniform distribution
$\bigl(\mathbf{A},\,\mathbf{b}\bigr) \gets \Rq^{k \times \ell} \times \Rq^k$.
The advantage of an adversary $\mathcal{B}$ is
\[
\Adv^{\MLWE}_{q, k, \ell, \eta}(\mathcal{B})
=
\bigl|\Pr[\mathcal{B}(\mathbf{A}, \mathbf{A}\mathbf{s}_1+\mathbf{s}_2) = 1]
- \Pr[\mathcal{B}(\mathbf{A}, \mathbf{b}) = 1]\bigr|.
\]
\end{definition}

The standard search-to-decision reduction (Langlois--Stehl\'e
\cite{lattice-estimator} cite-tree from FIPS~204~\S2 references) gives
the corresponding search hardness. The parameter sets in
\S\ref{sec:parameters} fix the concrete instances at NIST PQ
categories $2$, $3$, $5$.

\subsection{Pedersen commitments over $\Rq^k$}
\label{sec:prelim-pedersen}

Pulsar's DKG and reshare protocols build on a module-shaped
Pedersen commitment. The construction is the natural lift of the
$\Rq$ Pedersen commitment used by the Pulsar reference
implementation~\cite{boschini2024corona} to the module
ring $\Rq^k$.

\begin{construction}[Module Pedersen commitment]
\label{con:pedersen}
Let $\mathbf{A}, \mathbf{B} \in \Rq^{k \times \ell}$ be two independent
uniformly distributed public matrices. The Pedersen commitment to a
secret vector $\mathbf{c} \in \Rq^\ell$ with binding randomness
$\mathbf{r} \in \Rq^\ell$ is
\[
\Comm_{\mathbf{A}, \mathbf{B}}(\mathbf{c}; \mathbf{r}) \;=\;
\mathbf{A} \cdot \mathbf{c} + \mathbf{B} \cdot \mathbf{r} \pmod q.
\]
$\mathbf{r}$ is drawn from the FIPS~204 centred binomial $\chi_\eta^\ell$ in
the DKG instantiation. The commitment is $k$-vector valued in
$\Rq^k$.
\end{construction}

\begin{theorem}[Computational hiding]
\label{thm:pedersen-hiding}
Under $\MLWE_{q, k, \ell, \eta}$ in the matrix $\mathbf{B}$,
$\Comm_{\mathbf{A}, \mathbf{B}}$ is computationally hiding. Formally,
for any $\PPT$ adversary $\mathcal{A}$ choosing
$\mathbf{c}_0, \mathbf{c}_1 \in \Rq^\ell$,
\[
\Adv^{\mathrm{hide}}_{\mathbf{A}, \mathbf{B}}(\mathcal{A})
\le
\Adv^{\MLWE}_{q, k, \ell, \eta}(\mathcal{B})
\]
for an efficient reduction $\mathcal{B}$.
\end{theorem}

\begin{proof}[Proof sketch]
Identical to the BDLOP18~\cite{boschini2024corona} commitment-hiding
argument applied per Shamir-coordinate slot of the module ring.
Given an M-LWE instance $(\mathbf{B}, \mathbf{B}\mathbf{r} + \mathbf{e})$
vs.\ $(\mathbf{B}, \mathbf{u})$, the reduction sets
$\Comm_{\mathbf{A}, \mathbf{B}}(\mathbf{c}_b; \mathbf{r})
= \mathbf{A}\mathbf{c}_b + (\mathbf{B}\mathbf{r} + \mathbf{e})$, embeds
the bit $b$ into the choice of plaintext, and uses $\mathcal{A}$'s
guess. The noise term $\mathbf{e}$ is added back into the commitment
because the FIPS~204 noise distribution is fixed and the simulator
must produce exactly that distribution. The advantage
relationship is tight. \qed
\end{proof}

\begin{theorem}[Computational binding]
\label{thm:pedersen-binding}
Under $\MSIS_{q, k, 2\ell, \beta_{\mathrm{bind}}}$ in the wide
concatenation $[\mathbf{A} \,|\, \mathbf{B}]$, the Pedersen commitment
is computationally binding. The binding norm bound
$\beta_{\mathrm{bind}}$ is bounded by the maximum norm of any
honestly-produced opening:
$\beta_{\mathrm{bind}} \ge 2\eta\sqrt{n \cdot \ell}$.
\end{theorem}

\begin{proof}[Proof sketch]
Suppose $\mathcal{A}$ produces
$(\mathbf{c}_0, \mathbf{r}_0) \ne (\mathbf{c}_1, \mathbf{r}_1)$ with
$\Comm_{\mathbf{A},\mathbf{B}}(\mathbf{c}_0;\mathbf{r}_0)
= \Comm_{\mathbf{A},\mathbf{B}}(\mathbf{c}_1;\mathbf{r}_1)$. Then
$[\mathbf{A}\,|\,\mathbf{B}] \cdot
(\mathbf{c}_0 - \mathbf{c}_1,\,\mathbf{r}_0 - \mathbf{r}_1)^\top = 0$,
producing a short solution to MSIS in the concatenation, of norm
at most $2 \cdot \max(\|\mathbf{c}\|_\infty, \|\mathbf{r}\|_\infty)$,
which is bounded under the openings of interest. \qed
\end{proof}

\paragraph{Public matrix derivation.}
$\mathbf{A}$ is derived exactly as in FIPS~204 via
\texttt{ExpandA} (FIPS~204~\S7.1): \(\mathbf{A} = \texttt{ExpandA}(\rho)\),
$\rho \in \{0,1\}^{256}$ a public seed. $\mathbf{B}$ is derived from
a separate nothing-up-my-sleeve seed
$\rho_B = \cSHAKE(\text{empty}, 256, \text{`Pulsar'},
\text{`PULSAR-EXPANDB-V1'}\,)$. The two derivations are
domain-separated so that $\mathbf{A}$ and $\mathbf{B}$ are independent
even with shared randomness, satisfying the independence
requirement of \Cref{thm:pedersen-hiding}.

\subsection{Shamir secret sharing over $\Rq^k$}
\label{sec:prelim-shamir}

Shamir secret sharing extends to $\Rq^k$ slot-wise: each $\Zq$
coefficient of each polynomial of each module entry is shared
independently. Concretely, for a secret $\mathbf{s} \in \Rq^\ell$,
the dealer chooses a vector polynomial
\[
P(X) = \mathbf{s} + \mathbf{a}_1 X + \cdots + \mathbf{a}_{t-1} X^{t-1}
\in \bigl(\Rq^\ell\bigr)[X], \qquad \mathbf{a}_d \gets \Rq^\ell,
\]
and assigns party $i \in \{1, \dots, n\}$ the share $\mathbf{s}_i = P(i)$.
Reconstruction at $X = 0$ uses standard Lagrange interpolation in
$\Zq$ per coefficient slot:
\[
\mathbf{s} = \sum_{i \in T} \lambda_i^T \cdot \mathbf{s}_i, \qquad
\lambda_i^T = \prod_{j \in T,\, j \ne i} \frac{-j}{i - j} \pmod q,
\]
for any threshold subset $T \subseteq \{1, \dots, n\}$ with
$|T| = t$. Because the Lagrange coefficients $\lambda_i^T$ live in
$\Zq$, they act on $\mathbf{s}_i \in \Rq^\ell$ as scalar
multiplication --- the ring structure is preserved trivially. The
information-theoretic security of $(t-1)$-Shamir over $\Zq$ gives
the corresponding security over $\Rq^k$ per slot.

\paragraph{Dual sharing $(\mathbf{s}_i, \mathbf{u}_i)$.}
Pulsar's DKG produces a Shamir share of the FIPS~204 secret
$\mathbf{s}_1$ alongside a Shamir share of the Pedersen blinding
randomness $\mathbf{u}$. We denote the per-party share pair by
$(\mathbf{s}_i, \mathbf{u}_i)$. Both are shared at the same threshold
$t$ over the same evaluation-point set. The pair is bound to the
party's identity by the Pedersen commit verification at DKG Round~2.

\subsection{Discrete Gaussian sampling}
\label{sec:prelim-gauss}

Pulsar's signing protocol re-uses every distribution from FIPS~204
verbatim: the mask $\mathbf{y} \gets U_{\gamma_1}^\ell$ (uniform interval),
the secret centred binomial $\chi_\eta$, the challenge ternary
$B_\tau$. The Pedersen DKG additionally uses the centred binomial
$\chi_\eta$ for the blinding factor $\mathbf{r}$ to match
FIPS~204's secret distribution.

\paragraph{Why no Gaussian in $\Sign$.}
FIPS~204 uses uniform-interval masking, not Gaussian masking (this is
the Dilithium variant, not the rounded-Gaussian variant of
Falcon-like schemes). To produce byte-equal FIPS~204 signatures,
Pulsar must produce a $z$ vector with the same distribution as
\texttt{ML-DSA.Sign}'s $z$ vector --- which means uniform-interval
masking. We do \emph{not} use the Gaussian sampler of FIPS~204 \S3.3
in the threshold $\Sign$ path. We \emph{do} use the centred binomial
for the DKG blinding randomness $\mathbf{r}$ inside Pedersen
commitments; this distribution choice is internal to the DKG and is
not visible in the output FIPS~204 signature.

\paragraph{Constant-time samplers.}
All samplers used in Pulsar must be constant-time in their inputs.
The FIPS~204 reference rejection samplers are constant-time provided
the bound check is performed with branchless arithmetic; this is
discussed further in \S\ref{sec:impl-ct}.

\subsection{Rejection sampling in FIPS~204}
\label{sec:prelim-rejection}

FIPS~204~\S6.1 specifies the rejection-with-restart structure of
ML-DSA signing. The signer samples a fresh $\mathbf{y}$ each
iteration, computes $\mathbf{w} = \mathbf{A}\mathbf{y}$, derives the
challenge $c$ from the high-bits of $\mathbf{w}$ and the message
hash $\mu$, computes the candidate response
$\mathbf{z} = \mathbf{y} + c \cdot \mathbf{s}_1$, and either accepts
or restarts depending on two norm checks: $\|\mathbf{z}\|_\infty$
must lie in the bounded interval
$[-\gamma_1 + \beta + 1, \gamma_1 - \beta - 1]$, and
$\|\LowBits(\mathbf{w} - c\mathbf{s}_2, 2\gamma_2)\|_\infty$ must lie
in $[-\gamma_2 + \beta + 1, \gamma_2 - \beta - 1]$. On a fail, the
signer restarts with a fresh $\mathbf{y}$.

This restart structure interacts subtly with threshold signing: each
restart consumes fresh masking randomness, but the per-party masking
randomness must be domain-separated across restarts so that an
attacker observing multiple restarts on the same $(\mathit{pk}, \mu)$
learns nothing more than from a single attempt. We address this in
\S\ref{sec:impl-rejection-restart}.

% =====================================================================
\section{Algorithms}
\label{sec:algorithms}
% =====================================================================

\subsection{Distributed key generation \texorpdfstring{$\DKG$}{DKG}}
\label{sec:dkg}

The DKG produces a group public key $\mathit{pk}$ that is a fresh
FIPS~204 ML-DSA public key together with per-party threshold shares
of the master 32-byte ML-DSA seed.

\paragraph{Two instantiation paths.}
Pulsar admits two protocol families that produce the same FIPS~204
output shape:
\begin{description}[leftmargin=*, itemsep=0pt]
\item[v0.2 target (Pedersen-VSS, this section).] A degree-$(t-1)$
Shamir polynomial over $\Rq^\ell$ with Pedersen polynomial commits
in $\Rq^k$. The aggregated commit $\mathbf{t}$ identifies with the
FIPS~204 public-key inner vector $\mathbf{A}\mathbf{s}_1 +
\mathbf{B}\mathbf{u}$, and the substitution $\mathbf{s}_2 =
\mathbf{B}\mathbf{u}$ is M-LWE-indistinguishable from a fresh
$\chi_\eta^k$ sample (\Cref{thm:outind}). No party ever reconstructs
the master secret. The security analysis (\S\ref{sec:security}) is
stated on this protocol.
\item[v0.1 reference (Path~A reconstruction aggregator).] A
byte-level instantiation realised by the reference Go code
(\texttt{ref/go/pkg/pulsar}, v0.2.0). Each party samples a 32-byte
contribution $c_i$, byte-wise Shamir-shares $c_i$ over $\GF(257)$ to
the committee, and ML-KEM-768-wraps each per-recipient envelope
against the recipient's long-term identity public key. Round~3
reconstructs the master seed $\sigma = \cSHAKE(\sum_i c_i \,\|\,
\mathit{root})$ in every committee member's memory and runs
\texttt{ML-DSA.KeyGen} verbatim on $\sigma$. Threshold sign
reconstructs $\sigma$ at sign time and runs \texttt{ML-DSA.Sign}
verbatim, which is what makes the byte-equality argument
(\Cref{thm:outind}) trivial under this path. The trust model is
documented in \S\ref{sec:knownlimits}.
\end{description}
The remainder of this subsection specifies the v0.2 protocol;
``v0.1 note'' paragraphs flag the deviations realised by the reference
implementation.

\paragraph{Setup (out of band).}
Before $\DKG$, all parties agree on the common public parameters
$(q, n_{\mathrm{ring}}, k, \ell, \eta, \tau, \gamma_1, \gamma_2, \omega, \beta, d, \lambda)$
matching the chosen FIPS~204 parameter set (one of $\{44, 65, 87\}$),
plus the threshold pair $(n, t)$ with $n \ge 2t-1$.

\paragraph{Setup: identity directory (v0.1 reference, CR-7 + CR-8
closure).}
Each party publishes a long-term identity public key
$\mathit{IPK}_i = (\mathit{KEMpk}_i, \mathit{DSApk}_i)$ pairing an
ML-KEM-768 encapsulation key with an ML-DSA-65 verification key,
typically anchored in the chain's validator-set record. Every party
constructs an $\textsc{IdentityDirectory}: \mathit{NodeID} \mapsto
\mathit{IPK}$ covering every committee member before DKG begins.
Round-1 envelopes are KEM-sealed against each recipient's
$\mathit{KEMpk}$; threshold-sign per-pair MAC keys are derived from
ephemeral ML-KEM-768 session secrets authenticated under the long-term
ML-DSA-65 keys (see \S\ref{sec:sign}).

\paragraph{Public matrices.}
$\mathbf{A} = \texttt{ExpandA}(\rho)$ from a public seed $\rho$
(typically the chain genesis seed bound to the validator-set epoch),
and $\mathbf{B}$ from the domain-separated derivation in
\S\ref{sec:prelim-pedersen}.

\paragraph{Round 1 (commit and deal).}
Each party $i \in \{1, \dots, n\}$:
\begin{enumerate}[leftmargin=*, itemsep=0pt]
\item Samples a degree-$(t-1)$ polynomial in $X$ over $\Rq^\ell$:
\[
f_i(X) = \mathbf{c}_{i,0} + \mathbf{c}_{i,1} X + \cdots + \mathbf{c}_{i,t-1} X^{t-1},
\]
with $\mathbf{c}_{i,k} \gets \chi_\eta^\ell$ for each $k$. ($\mathbf{c}_{i,0}$ is party $i$'s contribution to the master secret $\mathbf{s}_1$.)
\item Samples a degree-$(t-1)$ polynomial in $X$ over $\Rq^\ell$:
\[
g_i(X) = \mathbf{r}_{i,0} + \mathbf{r}_{i,1} X + \cdots + \mathbf{r}_{i,t-1} X^{t-1},
\]
with $\mathbf{r}_{i,k} \gets \chi_\eta^\ell$ for each $k$.
\item Computes the Pedersen commitment to each coefficient pair:
\[
\mathbf{C}_{i,k} = \mathbf{A} \cdot \mathbf{c}_{i,k} + \mathbf{B} \cdot \mathbf{r}_{i,k}
\in \Rq^k, \quad k \in \{0, 1, \dots, t-1\}.
\]
\item Broadcasts $\{\mathbf{C}_{i,k}\}_{k=0}^{t-1}$ to all parties.
\item Privately sends $(\text{share}_{i \to j}, \text{blind}_{i \to j}) = (f_i(j), g_i(j))$ to each party $j \ne i$ over an authenticated point-to-point channel.
\end{enumerate}

\paragraph{Round 1.5 (cross-party commitment consistency).}
Each party broadcasts the commit digest
\[
H_i = \cSHAKE(\,\text{empty},\,256,\,\text{`Pulsar'},\,
\text{`PULSAR-DKG-COMMIT-V1'}\,||\,\mathit{encode}(\{\mathbf{C}_{i,k}\}_{k=0}^{t-1})\,).
\]
Parties compare digests received from every sender across the cohort.
Mismatch $\Rightarrow$ equivocation: the sender broadcast distinct
commitment vectors to distinct recipients. A signed
\textsc{ComplaintEquivocation} disqualifies the sender. This step is
load-bearing for identifiable abort.

\paragraph{Round 2 (verification and aggregation).}
Each party $j$ verifies, for every sender $i$:
\[
\mathbf{A} \cdot f_i(j) + \mathbf{B} \cdot g_i(j) \;\stackrel{?}{=}\;
\sum_{k=0}^{t-1} j^k \cdot \mathbf{C}_{i,k} \pmod q.
\]
The check is exact (both sides are $\Zq$-linear) and is performed
with a constant-time module-level comparison (no early exit on the
first mismatched coefficient slot --- see
\S\ref{sec:impl-ct}). Mismatch $\Rightarrow$ a signed
\textsc{ComplaintBadDelivery} naming sender $i$ and carrying the
(share, blind, commits) as evidence.

On success, party $j$ aggregates:
\[
\mathbf{s}_j = \sum_{i=1}^n f_i(j), \qquad
\mathbf{u}_j = \sum_{i=1}^n g_i(j), \qquad
\mathbf{t} = \sum_{i=1}^n \mathbf{C}_{i,0}.
\]
$\mathbf{s}_j$ and $\mathbf{u}_j$ are party $j$'s share of $\mathbf{s}_1$
and the Pedersen blinding $\mathbf{u}$, respectively. The aggregated
Pedersen commitment $\mathbf{t} \in \Rq^k$ identifies with the FIPS~204
public-key inner vector by construction:
\[
\mathbf{t} = \mathbf{A} \cdot \mathbf{s}_1 + \mathbf{B} \cdot \mathbf{u}
\;\equiv_{\mathrm{dist}}\;
\mathbf{A} \cdot \mathbf{s}_1 + \mathbf{s}_2,
\]
where the substitution $\mathbf{s}_2 = \mathbf{B}\mathbf{u}$ is
computationally indistinguishable from a fresh
$\chi_\eta^k$ sample under $\MLWE_{q,k,\ell,\eta}$ in $\mathbf{B}$
(\Cref{thm:outind}).

\paragraph{Output: FIPS~204-shaped $\mathit{pk}$.}
The group public key follows the FIPS~204 layout
(FIPS~204~\S5.1, Algorithm~5):
\[
\mathit{pk} = \bigl(\rho,\,\mathbf{t}_1\bigr), \qquad
(\mathbf{t}_1, \mathbf{t}_0) = \Power(\mathbf{t}, d).
\]
The high-bits component $\mathbf{t}_1$ is in $\mathit{pk}$; the low-bits
$\mathbf{t}_0$ is required by every signer at Round-2 aggregation
(Algorithm~\ref{alg:sign-agg} line~3 uses it in the FIPS~204 hint
formula) and is therefore stored alongside the aggregated public key
by every party. $\mathbf{t}_0$ does \emph{not} leak the secret
$\mathbf{s}_1$: it is the low-bits of a uniform-modulus quantity and
behaves as a public protocol parameter.

The aggregated Pedersen blinding $\mathbf{u} = \sum_i \mathbf{r}_{i,0}$
is reconstructable by any threshold quorum via Lagrange interpolation
on the $\{\mathbf{u}_j\}_{j \in T}$. In the strict FIPS~204 layout
$\mathbf{s}_2$ is part of the private signing key; in Pulsar
$\mathbf{u}$ is Shamir-shared with the same threshold structure as
$\mathbf{s}_1$.

\begin{remark}[$\mathbf{s}_2$ as derived quantity]
The FIPS~204 specification treats $\mathbf{s}_1$ and $\mathbf{s}_2$
as independent samples from $\chi_\eta$ (FIPS~204~\S5.1). Pulsar's
construction derives $\mathbf{s}_2$ from the Pedersen blinding
$\mathbf{u}$ via the structural identity
$\mathbf{s}_2 = \mathbf{B} \cdot \mathbf{u}$. Because $\mathbf{B}$ is
uniformly distributed and $\mathbf{u}$ is from $\chi_\eta^\ell$,
$\mathbf{s}_2$ is computationally indistinguishable from a fresh
$\chi_\eta^k$ sample under $\MLWE_{q, k, \ell, \eta}$; this is the
content of the output-indistinguishability theorem
(\Cref{thm:outind} below).
\end{remark}

\paragraph{Identifiable abort.}
Round~1.5 and Round~2 each produce signed complaint records as a
side-effect of detected deviations. The signing key is the party's
long-term identity key (e.g.\ the Ed25519 key advertised in the
chain's validator-set record). Complaint aggregation rules are
specified per the system model (\S\ref{sec:system-model}).

\subsection{Threshold signing \texorpdfstring{$\Sign$}{Sign}}
\label{sec:sign}

The threshold $\Sign$ protocol takes:

\begin{itemize}[leftmargin=*, itemsep=0pt]
\item the group public key $\mathit{pk}$,
\item the quorum $T \subseteq \{1, \dots, n\}$ with $|T| = t$,
\item per-party shares $\{(\mathbf{s}_i, \mathbf{u}_i)\}_{i \in T}$,
\item a message $\mu$ (after FIPS~204 message hash binding $\mathrm{tr}$),
\item a session id $\mathit{sid}$ unique per signature attempt,
\end{itemize}

and produces a FIPS~204 ML-DSA signature
$\sigma = (\tilde c, \mathbf{z}, \mathbf{h})$.

The protocol runs in two rounds plus a possibly-iterated rejection-restart
outer loop. The inner rejection-restart loop is parameterised by an
attempt counter $\kappa$; each new attempt advances $\kappa$ and
re-domain-separates all sampler PRNG seeds (see
\S\ref{sec:impl-rejection-restart}).

\begin{algorithm}[H]
\caption{$\Sign$ Round~1 (party $i$, attempt $\kappa$)}
\label{alg:sign-r1}
\begin{algorithmic}[1]
\State \textbf{Input:} $\mathit{pk}$, $\mathit{sid}$, $\kappa$, share $\mathbf{s}_i$,
quorum $T$, party-pair PRF keys $\{K_{i,j}\}_{j \in T \setminus \{i\}}$,
message hash $\mu$.
\State $\tau_1 \gets (\mathit{sid}, \kappa, T, i, \mathit{pk}, \mu)$
\Comment{Round-1 transcript binding}
\State $\text{seed}_i \gets \KMAC(K_i^{\mathrm{prng}},\,\tau_1,\,256,\,\text{`PULSAR-SIGN-PRNG-V1'})$
\Comment{Per-round PRNG seed (\S\ref{sec:impl-rejection-restart})}
\State Sample $\mathbf{y}_i \gets U_{\gamma_1}^\ell$ from the PRNG keyed by $\text{seed}_i$
\State Compute commit $\mathbf{w}_i = \mathbf{A} \mathbf{y}_i \in \Rq^k$
\State Compute high-bits commit $\bar{\mathbf{w}}_i = \HighBits(\mathbf{w}_i,\,2\gamma_2)$
\State Compute commit digest $D_i = \cSHAKE(\bar{\mathbf{w}}_i,\,256,\,\text{`Pulsar'},\,\text{`PULSAR-SIGN-R1-V1'} \| \tau_1)$
\State \textbf{For each} $j \in T \setminus \{i\}$ \textbf{do}
$\;\mathrm{MAC}_{i \to j} \gets \KMAC(K_{i,j},\,D_i \| \tau_1,\,256,\,\text{`PULSAR-SIGN-R1-MAC-V1'})$
\State \textbf{Broadcast:} $(D_i, \{\mathrm{MAC}_{i \to j}\}_{j \in T \setminus \{i\}})$
\end{algorithmic}
\end{algorithm}

\begin{algorithm}[H]
\caption{$\Sign$ Round~2 (party $i$, attempt $\kappa$)}
\label{alg:sign-r2}
\begin{algorithmic}[1]
\State \textbf{Input:} the received $\{D_j\}_{j \in T}$ and corresponding MACs, plus all Round-1 state from party $i$.
\State Verify every MAC received from each $j$ against the locally-known $K_{i,j}$. Abort with $\textsc{ComplaintMACFailure}(j)$ on mismatch.
\State Compute the aggregated commit hash $\bar D = \cSHAKE(\text{empty},\,256,\,\text{`Pulsar'},\,\text{`PULSAR-SIGN-R2-V1'} \,\|\, D_1 \| D_2 \| \cdots \| D_n)$
\State Recover the aggregated $\bar{\mathbf{w}} = \HighBits(\sum_{j \in T} \mathbf{w}_j,\,2\gamma_2)$ by summing per-party Round-1 commits (each party transmits its raw $\bar{\mathbf{w}}_i$ alongside the commit digest in Round~2 for verification --- the digest's role is binding under MAC, not concealment).
\State Compute the FIPS~204 challenge $\tilde c \gets \SHAKE(\mu \,\|\, \bar{\mathbf{w}},\,2\lambda)$ and expand to $c \in \Rq$ via $c = \texttt{SampleInBall}(\tilde c)$.
\State Compute the Lagrange coefficient $\lambda_i^T \in \Zq$ for party $i$ in the quorum $T$.
\State Compute party $i$'s response contribution:
$\mathbf{z}_i = \mathbf{y}_i + c \cdot \lambda_i^T \cdot \mathbf{s}_i \pmod q$
\State Compute party $i$'s low-bits contribution to the hint vector:
$\mathbf{r}_i = c \cdot \lambda_i^T \cdot \mathbf{u}_i$ \Comment{used in $\mathbf{h}$ computation by the aggregator}
\State \textbf{Broadcast:} $(\bar{\mathbf{w}}_i, \mathbf{z}_i, \mathbf{r}_i)$ alongside an integrity hash bound to the Round-2 transcript.
\end{algorithmic}
\end{algorithm}

\begin{algorithm}[H]
\caption{$\Sign$ Aggregation and rejection (any honest party in $T$)}
\label{alg:sign-agg}
\begin{algorithmic}[1]
\State Compute aggregate $\mathbf{z} = \sum_{j \in T} \mathbf{z}_j$ and aggregate $c\mathbf{s}_2 = \sum_{j \in T} \mathbf{r}_j$. By Lagrange linearity, $c\mathbf{s}_2 = c \cdot \mathbf{B}\mathbf{u}$ where $\mathbf{u} = \sum_{j} \lambda_j^T \mathbf{u}_j$ is the master Pedersen blinding (\S\ref{sec:dkg}).
\State Compute the FIPS~204 high-bits delta: let $\mathbf{w} = \sum_{j \in T} \mathbf{w}_j$ and $\mathbf{w}' = \mathbf{w} - c\mathbf{s}_2$.
\State Compute the candidate hint $\mathbf{h} = \MakeHint(-\langle c\mathbf{t}_0 \rangle,\,\mathbf{w}' + \langle c\mathbf{t}_0 \rangle,\,2\gamma_2)$ following FIPS~204~\S6.1 line 16, where $\mathbf{t}_0$ is the low-bits component of the public Pedersen-recovered commit (carried in $\mathit{pk}$'s private companion held by every party).
\State Compute the FIPS~204 rejection-check predicates (FIPS~204~\S6.1.1):
\State \quad $R_1$: $\|\mathbf{z}\|_\infty \stackrel{?}{<} \gamma_1 - \beta$
\State \quad $R_2$: $\|\LowBits(\mathbf{w}',\,2\gamma_2)\|_\infty \stackrel{?}{<} \gamma_2 - \beta$
\State \quad $R_3$: $\|\langle c\mathbf{t}_0\rangle\|_\infty \stackrel{?}{<} \gamma_2$
\State \quad $R_4$: $|\mathbf{h}|_{\mathrm{count}} \stackrel{?}{\le} \omega$
\State Compute the combined accept bit $\mathrm{accept} = R_1 \wedge R_2 \wedge R_3 \wedge R_4$, evaluating all $R_i$ unconditionally (constant-time, \S\ref{sec:impl-ct}).
\If{$\neg\mathrm{accept}$}
\State Restart: increment $\kappa$ and re-run from Round~1 with new PRNG seeds (\S\ref{sec:impl-rejection-restart}). \textbf{Crucially}: no per-round state is preserved across the restart boundary other than $(\mathit{sid}, \kappa, \mathit{pk}, \mu, T)$ and the long-term shares.
\Else
\State Output $\sigma = (\tilde c, \mathbf{z}, \mathbf{h})$.
\EndIf
\end{algorithmic}
\end{algorithm}

\begin{theorem}[Correctness of $\Sign$]
\label{thm:sign-correct}
If all parties in $T$ are honest and the rejection conditions
accept, then the aggregated $\sigma = (\tilde c, \mathbf{z}, \mathbf{h})$
satisfies $\texttt{ML-DSA.Verify}(\mathit{pk}, \mu, \sigma) = \mathit{accept}$.
\end{theorem}

\begin{proof}
By Lagrange-interpolation linearity:
\[
\mathbf{z}
= \sum_{j \in T} \bigl(\mathbf{y}_j + c \cdot \lambda_j^T \cdot \mathbf{s}_j\bigr)
= \mathbf{y} + c \cdot \mathbf{s}_1
\]
where $\mathbf{y} = \sum_{j \in T} \mathbf{y}_j$ is the aggregate
mask (distributed as uniform interval over the per-coordinate sum,
and within the FIPS~204 acceptance interval by the rejection
check) and $\sum_{j \in T} \lambda_j^T \mathbf{s}_j = \mathbf{s}_1$
by Lagrange interpolation at $X = 0$. This is exactly the FIPS~204
single-signer $\mathbf{z}$. The challenge $\tilde c$ is the
FIPS~204 challenge computed on the same
$(\mu, \HighBits(\mathbf{A}\mathbf{y}, 2\gamma_2))$ input. The hint
$\mathbf{h}$ is the FIPS~204 hint computed against the same
$(\mathbf{w}, c\mathbf{s}_2, 2\gamma_2)$ inputs. Therefore the triple
$(\tilde c, \mathbf{z}, \mathbf{h})$ is identical to what
single-signer FIPS~204 would output on the same
$(\mathit{pk}, \mu, \mathbf{y})$. Verification follows from FIPS~204's
correctness. \qed
\end{proof}

\begin{remark}[Round-1 commit-digest binding]
The Round-1 commit digest $D_i$ is bound under MAC, not concealment:
its purpose is to commit each party to its
$\bar{\mathbf{w}}_i$ before learning the challenge $c$. The
plaintext $\bar{\mathbf{w}}_i$ is transmitted in Round~2 alongside
the response; recipients re-derive $D_i$ from the received plaintext
and check equality. This is the standard MAC-and-commit pattern
inherited from Pulsar's $\Sign$ in \texttt{pulsar/sign/sign.go}.
\end{remark}

\subsection{Verification \texorpdfstring{$\Verify$}{Verify}}
\label{sec:verify}

$\Verify$ is \texttt{ML-DSA.Verify} from FIPS~204~\S6.3, verbatim. No
threshold-specific changes; no per-threshold extension fields; no
side parameters. We quote FIPS~204 for completeness:

\begin{algorithm}[H]
\caption{$\Verify$ (\texttt{ML-DSA.Verify} per FIPS~204 \S6.3)}
\label{alg:verify}
\begin{algorithmic}[1]
\State \textbf{Input:} $\mathit{pk} = (\rho, \mathbf{t}_1)$, message $\mu$, signature $\sigma = (\tilde c, \mathbf{z}, \mathbf{h})$.
\State $\mathbf{A} \gets \texttt{ExpandA}(\rho)$
\State $\mu' \gets \SHAKE(\mathrm{tr} \| \mu, 512)$ where $\mathrm{tr}$ is the FIPS~204 public-key fingerprint
\State $c \gets \texttt{SampleInBall}(\tilde c)$
\State $\mathbf{w}' \gets \UseHint(\mathbf{h},\,\mathbf{A}\mathbf{z} - c \cdot 2^d \mathbf{t}_1,\,2\gamma_2)$
\State $\tilde c' \gets \SHAKE(\mu' \| \mathbf{w}', 2\lambda)$
\State \textbf{Output:}
$\mathit{accept}$ iff
$\|\mathbf{z}\|_\infty < \gamma_1 - \beta$
and
$\tilde c = \tilde c'$
and
$|\mathbf{h}|_{\mathrm{count}} \le \omega$.
\end{algorithmic}
\end{algorithm}

Note that line 3's $\mu'$ definition --- the public-key fingerprint
prefixing --- is exactly the FIPS~204 binding. Pulsar's $\Sign$
must compute its challenge over the \emph{same} $\mu'$ for the
verifier to accept; this is the $\mu$ used in
Algorithm~\ref{alg:sign-r2} step 5.

This is the entire Class~N point: any FIPS-validated ML-DSA
verifier (BoringSSL FIPS, AWS-LC, OpenSSL~3, liboqs) consumes
Pulsar signatures without modification.

\subsection{Proactive resharing \texorpdfstring{$\Reshare$}{Reshare}}
\label{sec:reshare}

$\Reshare$ rotates the share distribution from an old committee
$\mathcal{O}$ with threshold $t_{\mathrm{old}}$ to a new committee
$\mathcal{N}$ with threshold $t_{\mathrm{new}}$ without (i) reconstructing
the master secret $\mathbf{s}_1$ in any party's memory, and (ii)
changing the public key $\mathit{pk}$. The construction is HJKY97-style
(\cite{hjky97}) lifted to $\Rq^k$.

\paragraph{Beacon-randomised quorum selection (DD-006).}
A na\"ive deterministic quorum selection (e.g.\ smallest-ID
$t_{\mathrm{old}}$ subset of $\mathcal{O}$) lets a mobile adversary
\emph{game} the corrupt-set positioning: the adversary, knowing in
advance which parties will be quorum members at each reshare, can
schedule its $t-1$ corruption budget to maximally cover the quorum.
The HIP-0077 red-review finding F10 made this concrete for Pulsar's
reshare layer. Pulsar mitigates by deriving the quorum from the
chain's randomness beacon $\beta_e$ for epoch $e$:
\[
\mathcal{Q}_e
\;=\;
\bigl\{\text{the } t_{\mathrm{old}} \text{ smallest-ID members of }
\mathcal{O}\text{ in the permutation induced by }
\cSHAKE(\beta_e, 256\cdot|\mathcal{O}|, \text{`Pulsar'},
\text{`PULSAR-RESHARE-BEACON-V1'})\bigr\}.
\]
The beacon $\beta_e$ must be unpredictable in advance to the
adversary; the Lux consensus randomness layer (Quasar) provides this
unpredictability. With an unbiasable beacon, the adversary's optimal
strategy reduces to randomly guessing which corrupted parties will be
selected, giving negligible advantage over uniform.

\paragraph{Protocol.}
\begin{enumerate}[leftmargin=*, itemsep=0pt]
\item Compute $\mathcal{Q}_e \subseteq \mathcal{O}$ as above with
$|\mathcal{Q}_e| = t_{\mathrm{old}}$.
\item Compute Lagrange coefficients $\lambda_i^{\mathcal{Q}_e}$ at
$X = 0$ for each $i \in \mathcal{Q}_e$.
\item Each $i \in \mathcal{Q}_e$ samples a degree-$(t_{\mathrm{new}} - 1)$
polynomial in $X$ over $\Rq^\ell$:
\[
g_i(X) = (\lambda_i^{\mathcal{Q}_e} \cdot \mathbf{s}_i) + \mathbf{r}_{i,1} X + \cdots + \mathbf{r}_{i, t_{\mathrm{new}}-1} X^{t_{\mathrm{new}}-1},
\]
with $\mathbf{r}_{i,d} \gets \Rq^\ell$ uniform for $d \ge 1$, and
similarly an independent polynomial $h_i(X)$ with constant term
$\lambda_i^{\mathcal{Q}_e} \mathbf{u}_i$ for the blinding share.
\item Each $i$ commits to $g_i, h_i$ via Pedersen commitments
$\mathbf{C}^g_{i,k}, \mathbf{C}^h_{i,k}$, broadcasts them, and
privately delivers $(g_i(j), h_i(j))$ to each $j \in \mathcal{N}$.
\item Each new party $j$ verifies, for every $i \in \mathcal{Q}_e$:
$\mathbf{A} g_i(j) + \mathbf{B} h_i(j) \stackrel{?}{=}
\sum_k j^k \mathbf{C}^g_{i,k} + \sum_k j^k \mathbf{C}^h_{i,k}$, with
the same constant-time module-level comparison as in
\S\ref{sec:dkg}.
\item Each $j$ aggregates
$\mathbf{s}'_j = \sum_{i \in \mathcal{Q}_e} g_i(j)$ and
$\mathbf{u}'_j = \sum_{i \in \mathcal{Q}_e} h_i(j)$.
\item All new parties produce an \textsc{ActivationCertificate}: a
threshold signature of the reshare transcript hash under the
\emph{unchanged} $\mathit{pk}$, proving the new committee is
operational. The chain accepts the new epoch only on activation
certificate verification.
\end{enumerate}

\begin{theorem}[Public-key invariance]
\label{thm:reshare-pkinv}
Let $\mathbf{s}'_j$ and $\mathbf{u}'_j$ be the shares produced by
$\Reshare$. For every threshold subset
$T' \subseteq \mathcal{N}$ with $|T'| = t_{\mathrm{new}}$,
\[
\sum_{j \in T'} \lambda_j^{T'} \mathbf{s}'_j = \mathbf{s}_1 \quad\text{and}\quad
\sum_{j \in T'} \lambda_j^{T'} \mathbf{u}'_j = \mathbf{u} \pmod q.
\]
Therefore $\mathit{pk}$ is invariant under $\Reshare$.
\end{theorem}

\begin{proof}
Define $G(X) = \sum_{i \in \mathcal{Q}_e} g_i(X)$. $G$ has degree at
most $t_{\mathrm{new}} - 1$ and constant term
\[
G(0) = \sum_{i \in \mathcal{Q}_e} g_i(0)
= \sum_{i \in \mathcal{Q}_e} \lambda_i^{\mathcal{Q}_e} \mathbf{s}_i
= \mathbf{s}_1,
\]
where the last equality is Lagrange interpolation over the old
quorum. The new shares are $\mathbf{s}'_j = G(j)$ for $j \in \mathcal{N}$,
so $\{\mathbf{s}'_j\}_{j \in \mathcal{N}}$ are valid Shamir shares of
$G(0) = \mathbf{s}_1$ with threshold $t_{\mathrm{new}}$. The same argument
applies to $\{\mathbf{u}'_j\}_{j \in \mathcal{N}}$. \qed
\end{proof}

\paragraph{Same-committee refresh.}
We also specify a $\Refresh$ primitive (HJKY97 zero-polynomial form)
for periodic mobile-adversary hygiene within a stable validator set.
$\Refresh$ is structurally identical to $\Reshare$ with the
substitution $g_i(0) = 0, h_i(0) = 0$ (zero constant term), the same
committee on both sides, and an updated activation certificate. Its
correctness and security arguments are direct corollaries of
\Cref{thm:reshare-pkinv} (every honest party's polynomial vanishes at
$X = 0$, so the master secret is preserved).

\subsection{Identifiable abort}
\label{sec:abort}

Every detected deviation in $\DKG$, $\Sign$, or $\Reshare$ produces
a signed complaint. A complaint is a tuple
$(\text{kind}, \text{accuser}, \text{accused}, \text{epoch}, \text{evidence})$
where:

\begin{itemize}[leftmargin=*, itemsep=0pt]
\item \textsc{ComplaintEquivocation}: the accused party broadcast distinct commitment vectors to distinct recipients. Evidence: two commits and the signed broadcast from the accused.
\item \textsc{ComplaintBadDelivery}: the accused party's private share/blind to the accuser fails Pedersen verification. Evidence: (share, blind, commits) showing the verification fails.
\item \textsc{ComplaintMACFailure}: a MAC from the accused fails verification. Evidence: the failing MAC and the recipient's key.
\item \textsc{ComplaintRangeFailure}: the accused party's contribution would have caused the aggregated signature to fail the FIPS~204 norm check by an amount inconsistent with honest behaviour. Evidence: the per-party transcript line.
\end{itemize}

All complaints are signed under the accuser's long-term identity key
(e.g.\ Ed25519). The system model's slashing layer
(\S\ref{sec:system-model}) consumes complaints; a quorum of distinct
accusers against the same accused is a disqualification event.

% =====================================================================
\section{Security analysis}
\label{sec:security}
% =====================================================================

The formal game definitions are stated in the companion file
\texttt{spec/security-games.tex}, included below:

% Pulsar security game definitions --- included by pulsar.tex
% \section{Security analysis}
%
% Standard cryptographic-game style. Each game pins the challenger
% setup, the adversary's oracle interface, the winning condition, and
% the advantage; reductions are stated to the underlying assumption.

\subsection{Adversary model and notation}
\label{sec:sec-notation}

Throughout this section the adversary $\mathcal{A}$ is probabilistic
polynomial-time (PPT) and is parameterised by the corruption
threshold $t-1$ (the largest number of parties $\mathcal{A}$ may
corrupt). Honest parties' channels are authenticated; corrupted
parties act under $\mathcal{A}$'s control. We write $\mathcal{H}
\subseteq \{1, \dots, n\}$ for the honest-party set and
$\mathcal{C} = \{1, \dots, n\} \setminus \mathcal{H}$ for the corrupted
set; $|\mathcal{C}| \le t-1$.

The challenger maintains a transcript log
$\mathcal{T}$ recording every message visible to $\mathcal{A}$
(broadcast messages, point-to-point messages into corrupted parties,
public bulletin entries). The adversary's oracle interface depends on
the game; we specify it inline.

We write $\mathcal{H}_{\mathsf{RO}}$ for the random oracle used to model the
hash family $\{\SHAKE, \cSHAKE, \KMAC, \TupleHash\}$ in the ROM
analysis. Distinct customisation tags yield independent random
oracles by the indifferentiability argument of
SP~800-185~\cite{nist-sp800-185}.

% =====================================================================
\subsection{Game TS-UF (threshold strong unforgeability)}
\label{sec:game-tsuf}
% =====================================================================

\begin{game}[TS-UF]
\label{game:tsuf}
The challenger:
\begin{enumerate}[leftmargin=*, itemsep=0pt]
\item Runs $\DKG$ honestly across $n$ simulated parties; the adversary
sees every transcript message and learns the shares
$\{(\mathbf{s}_i, \mathbf{u}_i)\}_{i \in \mathcal{C}}$ of corrupted
parties. The group public key $\mathit{pk}$ is announced.
\item Exposes a $\Sign$ oracle $\mathcal{O}_{\Sign}(\mu, T)$: the
adversary chooses a message $\mu$ and a quorum
$T \subseteq \{1, \dots, n\}$ with $|T| = t$ and
$|T \cap \mathcal{C}| \le t-1$. The challenger simulates honest
parties' contributions to $\Sign(\mu, T)$ and returns the full
transcript including the final FIPS~204 signature $\sigma$.
\item Exposes a $\mathcal{H}_{\mathsf{RO}}$ oracle for all FIPS~202 / SP~800-185 hash
calls.
\end{enumerate}

The adversary outputs a forgery $(\mu^*, \sigma^*)$. $\mathcal{A}$
wins if:
\begin{enumerate}[label=(\roman*), leftmargin=*, itemsep=0pt]
\item $\texttt{ML-DSA.Verify}(\mathit{pk}, \mu^*, \sigma^*) = \mathit{accept}$,
\item $\mathcal{A}$ did not query $\mathcal{O}_{\Sign}(\mu^*, \cdot)$ with
the same $\mu^*$ producing $\sigma^*$, and
\item $|\mathcal{C}| \le t-1$ (the corruption budget was respected).
\end{enumerate}
$\mathcal{A}$'s advantage is
$\Adv^{\mathsf{TS\textsf{-}UF}}_{\mathrm{Pulsar\text{-}M}}(\mathcal{A}) = \Pr[\mathcal{A} \text{ wins}]$.
\end{game}

\begin{theorem}[TS-UF security]
\label{thm:tsuf}
For every $\PPT$ adversary $\mathcal{A}$ against \Cref{game:tsuf} with
$q_S$ signing queries and $q_H$ random-oracle queries, there exist
efficient reductions $\mathcal{B}_{\MLWE}, \mathcal{B}_{\MSIS},
\mathcal{B}_{\mathrm{ML\text{-}DSA}}$ such that
\[
\Adv^{\mathsf{TS\textsf{-}UF}}_{\mathrm{Pulsar\text{-}M}}(\mathcal{A})
\;\le\;
\Adv^{\MLWE}_{q, k, \ell, \eta}(\mathcal{B}_{\MLWE})
+ \Adv^{\MSIS}_{q, k, 2\ell, \beta_{\mathrm{bind}}}(\mathcal{B}_{\MSIS})
+ \Adv^{\mathsf{EUF\textsf{-}CMA}}_{\mathrm{ML\text{-}DSA}}(\mathcal{B}_{\mathrm{ML\text{-}DSA}})
+ \frac{q_S \cdot q_H}{2^\lambda}.
\]
\end{theorem}

\begin{proof}[Proof sketch]
We argue via a sequence of hybrids:

\paragraph{Hybrid $H_0$.} The real TS-UF experiment.

\paragraph{Hybrid $H_1$.} The challenger simulates the $\Sign$ oracle
without using the honest parties' shares. Specifically, the
simulator chooses the final signature $\sigma = (\tilde c, \mathbf{z},
\mathbf{h})$ by running \texttt{ML-DSA.Sign} against $\mathit{pk}$
\emph{as if it knew the secret key}, then programs the per-party
transcript so the aggregation lands on the chosen $\sigma$. The
honest-party Round-1 commits are set to $\mathbf{w}_j^{(\mathrm{sim})} =
\mathbf{A}\mathbf{y}_j^{(\mathrm{sim})}$ for simulated masks $\mathbf{y}_j^{(\mathrm{sim})}$
drawn uniformly under the $U_{\gamma_1}$ constraint and consistent with
the chosen $\mathbf{z}$ via
$\sum_j \mathbf{y}_j^{(\mathrm{sim})} = \mathbf{z} - c\mathbf{s}_1$.
Random-oracle programming sets the $\bar{\mathbf{w}}$-derived challenge to
the chosen $\tilde c$. The distinguishing advantage between $H_0$
and $H_1$ is bounded by $q_S \cdot q_H / 2^\lambda$ (probability of a
collision in the RO programming).

\paragraph{Hybrid $H_2$.} The challenger replaces the Pedersen commit
matrix $\mathbf{B}$'s structure --- $\mathbf{B}$ is now a fresh
uniform matrix \emph{independent} of $\mathbf{A}$. This is identical
to $H_1$ because $\mathbf{B}$ was already derived from a separate
domain-separated seed in the construction; the hybrid is a no-op
distributionally and is included only for proof clarity.

\paragraph{Hybrid $H_3$.} The challenger replaces the DKG-time
Pedersen blinding $\mathbf{u}$ with a uniform sample from $\Rq^\ell$
(rather than $\chi_\eta^\ell$). Distinguishing advantage is bounded
by $\Adv^{\MLWE}_{q, k, \ell, \eta}(\mathcal{B}_{\MLWE})$ via the
hiding property of the Pedersen commitment
(\Cref{thm:pedersen-hiding}).

\paragraph{Hybrid $H_4$ (reduction to ML-DSA EUF-CMA).} In $H_4$ the
simulator obtains $\mathit{pk}$ from a single-party ML-DSA EUF-CMA
challenger; it answers $\mathcal{O}_{\Sign}(\mu, T)$ queries by
relaying the signing query to the ML-DSA challenger and using $H_1$'s
simulator to expand the resulting $\sigma$ into a per-party
transcript. Any adversary that wins TS-UF against this hybrid
produces a forgery $(\mu^*, \sigma^*)$ that wins the ML-DSA EUF-CMA
game (since the Pulsar $\Verify$ is unmodified ML-DSA
$\Verify$). Distinguishing advantage between $H_3$ and $H_4$ is
absorbed into the MSIS binding term: any deviation by the simulator
from the honest-party transcript shape is detected by binding of the
Pedersen commitment, contributing
$\Adv^{\MSIS}_{q, k, 2\ell, \beta_{\mathrm{bind}}}(\mathcal{B}_{\MSIS})$.

Summing the hybrid advantages gives the bound. \qed
\end{proof}

\begin{remark}[Tightness]
The reduction is tight modulo the standard $q_S q_H / 2^\lambda$
random-oracle programming penalty. For ML-DSA-65 ($\lambda = 192$),
this term is negligible for any practical $(q_S, q_H)$. The MSIS
term governs the Pedersen-binding side of the argument; the MLWE
term governs the hiding side; the ML-DSA term governs the underlying
single-party unforgeability. Pulsar inherits ML-DSA's EUF-CMA
bound verbatim modulo these threshold-protocol-specific terms.
\end{remark}

% =====================================================================
\subsection{Game OUT-IND (output indistinguishability from FIPS~204)}
\label{sec:game-outind}
% =====================================================================

\begin{game}[OUT-IND]
\label{game:outind}
The challenger flips a fair bit $b \in \{0, 1\}$ at setup and exposes
two oracles to $\mathcal{A}$:
\begin{itemize}[leftmargin=*, itemsep=0pt]
\item If $b = 0$: $\DKG$ produces $\mathit{pk}$ honestly, and
$\mathcal{O}_{\Sign}(\mu, T)$ runs Pulsar $\Sign$ honestly and
returns only the final $\sigma$ (transcript is hidden).
\item If $b = 1$: $\mathit{pk}$ is generated by single-party
\texttt{ML-DSA.KeyGen}, and $\mathcal{O}_{\Sign}(\mu, T)$ ignores the
quorum $T$ and runs single-party \texttt{ML-DSA.Sign}, returning
$\sigma$.
\end{itemize}
$\mathcal{A}$ outputs a guess $\hat b$. Advantage:
$\Adv^{\mathsf{OUT\textsf{-}IND}}_{\mathrm{Pulsar\text{-}M}}(\mathcal{A})
= |\Pr[\hat b = b] - 1/2|$.
\end{game}

\begin{theorem}[Output indistinguishability]
\label{thm:outind}
For every $\PPT$ adversary $\mathcal{A}$ against \Cref{game:outind},
\[
\Adv^{\mathsf{OUT\textsf{-}IND}}_{\mathrm{Pulsar\text{-}M}}(\mathcal{A})
\;\le\;
\Adv^{\MLWE}_{q, k, \ell, \eta}(\mathcal{B}_{\MLWE}) +
\frac{1}{2^{2\lambda}}.
\]
\end{theorem}

\begin{proof}[Proof sketch]
\textbf{$\mathit{pk}$ indistinguishability.}
A FIPS~204 public key is $\mathit{pk}_{\mathrm{F}} = (\rho, \Power(\mathbf{A}\mathbf{s}_1 + \mathbf{s}_2, d))$ with $\mathbf{s}_1, \mathbf{s}_2 \gets \chi_\eta$.
A Pulsar public key is $\mathit{pk}_{\mathrm{P}} = (\rho, \Power(\mathbf{A}\mathbf{s}_1 + \mathbf{B}\mathbf{u}, d))$ with $\mathbf{s}_1, \mathbf{u} \gets \chi_\eta$ (centred binomial) and $\mathbf{B}$ uniform.
The two distributions are identical except in the $\mathbf{s}_2$ vs.\ $\mathbf{B}\mathbf{u}$ component. Under $\MLWE_{q, k, \ell, \eta}$ in $\mathbf{B}$, $\mathbf{B}\mathbf{u}$ is computationally indistinguishable from a uniform $\mathbf{w} \in \Rq^k$; conditioned on uniformity, the post-$\Power$ projection is uniform over the range of $\Power$. Single-party ML-DSA's $\mathbf{s}_2 \sim \chi_\eta^k$ also projects through $\Power$ to a distribution that is statistically close to uniform over the $\Power$ image. The total distance between the two $\mathit{pk}$ distributions is therefore bounded by $\Adv^{\MLWE}_{q, k, \ell, \eta}(\mathcal{B}_{\MLWE})$ plus a negligible statistical term from the $\Power$ rounding.

\textbf{$\sigma$ indistinguishability.}
By correctness (\Cref{thm:sign-correct}), the Pulsar aggregated
$\sigma$ is the FIPS~204 signature on $(\mathit{pk}, \mu, \mathbf{y})$ where $\mathbf{y} = \sum_j \mathbf{y}_j$ is the aggregated mask. The marginal distribution of $\mathbf{y}$ (sum of $t$ uniform-interval samples) is a discrete convolution; FIPS~204's single-party mask is a single uniform-interval sample. The two distributions differ in their concentration profile but are projected through the rejection-sampling acceptance condition $\|\mathbf{z}\|_\infty < \gamma_1 - \beta$, which forces the accepted distribution into the same FIPS~204 acceptance region. Conditioned on acceptance, the per-coordinate distribution of $\mathbf{z}$ is uniform on $[-\gamma_1 + \beta + 1, \gamma_1 - \beta - 1]$ in both cases (this is the FIPS~204 rejection-sampling invariant). Therefore the post-rejection distribution is identical between Pulsar and FIPS~204, and the output $\sigma$ is distributionally identical given the same $\tilde c$ and $\mathbf{h}$.

The statistical term $1/2^{2\lambda}$ accounts for the negligible
probability of a $\tilde c$ collision in the RO. \qed
\end{proof}

\begin{remark}[Caveat on signature byte equality]
\Cref{thm:outind} states distributional equivalence. For \emph{byte
equality} on a specific signature, the Pulsar $\Sign$
implementation MUST emit exactly the FIPS~204 byte encoding of
$(\tilde c, \mathbf{z}, \mathbf{h})$ with the same canonicalisation
(coefficient lifting, sign convention, packing). This is an
implementation-level requirement, not a cryptographic theorem; the
encoding section (\S\ref{sec:encodings}) pins it. The byte-equality
property is what unlocks Class~N1.
\end{remark}

% =====================================================================
\subsection{Game ROB (robustness / identifiable abort)}
\label{sec:game-rob}
% =====================================================================

\begin{game}[ROB]
\label{game:rob}
$\mathcal{A}$ corrupts up to $t-1$ parties and chooses one of the
following deviation strategies:
\begin{enumerate}[label=(\alph*), leftmargin=*, itemsep=0pt]
\item Send distinct DKG commit vectors to distinct recipients (equivocation).
\item Send a private DKG share that fails Pedersen verification.
\item Send a Round-1 sign commit with an invalid MAC.
\item Send a Round-2 sign response that, combined with honest contributions, would push the aggregate $\mathbf{z}$ outside the FIPS~204 norm.
\item Refuse to participate at any round boundary.
\end{enumerate}
$\mathcal{A}$ wins if any honest party either (i) completes the
protocol against the deviation \emph{without} producing a valid
complaint identifying a corrupted deviator, or (ii) the honest
party's resulting state is inconsistent with at least one other
honest party's state (no recovery).
\end{game}

\begin{theorem}[Robustness]
\label{thm:rob}
For every $\PPT$ adversary $\mathcal{A}$ against \Cref{game:rob}, the
honest parties either produce a valid \textsc{Complaint}
record identifying a corrupted deviator, or successfully complete
the protocol. Formally,
$\Adv^{\mathsf{ROB}}_{\mathrm{Pulsar\text{-}M}}(\mathcal{A})$ is
bounded by the signature-forgery advantage on the parties' long-term
identity-key scheme (e.g.\ Ed25519) plus
$\Adv^{\MSIS}_{q, k, 2\ell, \beta_{\mathrm{bind}}}(\mathcal{B}_{\MSIS})$ for
deviation~(a),
plus the random-oracle collision term $q_H / 2^{256}$ for the
commit-digest binding.
\end{theorem}

\begin{proof}[Sketch]
For deviation (a), the Round-1.5 commit digest broadcast by every
honest recipient binds each sender to one commit vector across the
cohort (under collision-resistance of $\cSHAKE$). Two distinct
commits would produce two distinct digests, observable to every
recipient as a digest mismatch; the recipient signs a
\textsc{ComplaintEquivocation} carrying the two signed broadcasts as
evidence. The signature soundness of the accused party's long-term
identity key prevents the accused from repudiating its own
broadcasts; forging two broadcasts requires breaking the identity
signature.

For deviation (b), the Pedersen identity check
(\Cref{thm:pedersen-binding}) fails with probability $1 - \negl$ if
$\mathcal{A}$ deviates from a valid (share, blind, commit) triple.
The failing recipient signs a \textsc{ComplaintBadDelivery}
carrying the triple as evidence; the complaint is verifiable by any
third party because the public Pedersen check is determined by
$(\mathbf{A}, \mathbf{B})$.

For deviations (c), (d), (e): MAC verification, Round-2 norm
verification (against the public broadcast of $\bar{\mathbf{w}}_i,
\mathbf{z}_i, \mathbf{r}_i$), and round-window timeouts respectively
produce \textsc{ComplaintMACFailure}, \textsc{ComplaintRangeFailure},
\textsc{ComplaintTimeout} records, signed by the detecting party. Each
complaint is publicly verifiable from the transcript and the
accuser's identity-key signature. \qed
\end{proof}

% =====================================================================
\subsection{Game ADAPT (adaptive corruption)}
\label{sec:game-adapt}
% =====================================================================

\begin{game}[ADAPT]
\label{game:adapt}
Identical to \Cref{game:tsuf} except $\mathcal{A}$ may issue
\emph{adaptive} corruption queries: at any point during the protocol
execution, $\mathcal{A}$ chooses a previously-honest party $i$ to
corrupt; the challenger reveals party $i$'s entire current state
including secret share, current per-round randomness, and PRNG
seeds. The total number of corrupted parties at any time is bounded
by $t-1$.
\end{game}

\begin{theorem}[Adaptive corruption --- sketch only]
\label{thm:adapt-sketch}
\emph{Statement deferred to round-2 MPTC.} A formal
adaptive-corruption proof requires either (i) a complexity-leveraging
argument that loses a $\binom{n}{t-1}$ factor in the security
parameter and is therefore loose for practical $(n, t)$, or (ii) a
trapdoor-based simulation argument compatible with the Pedersen
commit structure. \cite{boschini2024corona} sketches direction (i) for the
Ring-LWE setting; the module-shaped adaptation is structural but
non-trivial.
\end{theorem}

\begin{openq}[Adaptive corruption]
Tight adaptive corruption for module-shaped two-round
threshold lattice signatures is an open research question. Pulsar
defers a formal proof to round 2; under the static-corruption model
\Cref{thm:tsuf} suffices for round 1.
\end{openq}

% =====================================================================
\subsection{Game MOB (mobile-adversary security, proactive resharing)}
\label{sec:game-mob}
% =====================================================================

\begin{game}[MOB]
\label{game:mob}
Multi-epoch experiment with epoch counter $e \in \{1, 2, \dots, E\}$.
At each epoch $e$:
\begin{enumerate}[leftmargin=*, itemsep=0pt]
\item $\mathcal{A}$ chooses a corruption set $\mathcal{C}_e \subseteq \{1, \dots, n\}$ with $|\mathcal{C}_e| \le t-1$. Corruptions are independent across epochs: a party in $\mathcal{C}_e \cap \mathcal{C}_{e+1}$ is corrupted across both epochs, but a party in $\mathcal{C}_e \setminus \mathcal{C}_{e+1}$ \emph{uncorrupts} between epochs (forgets its epoch-$e$ secrets but retains its identity-key state).
\item The challenger reveals the epoch-$e$ shares and per-round randomness of every party in $\mathcal{C}_e$.
\item $\mathcal{A}$ may issue $\mathcal{O}_{\Sign}$ queries within the epoch.
\item At the epoch boundary, the challenger runs $\Reshare$ honestly with a fresh beacon $\beta_{e+1}$. The beacon is unbiasable: $\mathcal{A}$ sees $\beta_{e+1}$ \emph{after} committing to $\mathcal{C}_{e+1}$.
\end{enumerate}
At the end of epoch $E$, $\mathcal{A}$ outputs a forgery
$(\mu^*, \sigma^*)$. $\mathcal{A}$ wins under the TS-UF winning
condition.
\end{game}

\begin{theorem}[Mobile-adversary security]
\label{thm:mob}
For every $\PPT$ adversary $\mathcal{A}$ against \Cref{game:mob} over
$E$ epochs,
\[
\Adv^{\mathsf{MOB}}_{\mathrm{Pulsar\text{-}M}}(\mathcal{A})
\;\le\;
E \cdot \Adv^{\mathsf{TS\textsf{-}UF}}_{\mathrm{Pulsar\text{-}M}}(\mathcal{B})
+ E \cdot \Adv^{\mathsf{beacon\textsf{-}unbias}}_{\mathrm{chain}}(\mathcal{B}_{\mathrm{beacon}}) +
E \cdot \Adv^{\MLWE}_{q, k, \ell, \eta}(\mathcal{B}_{\MLWE}).
\]
\end{theorem}

\begin{proof}[Proof sketch]
The proof proceeds by induction on epochs. The base case $e = 1$ is
TS-UF directly. For the inductive step, assume the master secret
$\mathbf{s}_1$ is hidden from $\mathcal{A}$ through epoch $e$. At the
epoch boundary:
\begin{enumerate}[leftmargin=*, itemsep=0pt]
\item The reshare quorum $\mathcal{Q}_{e+1}$ is selected from the
beacon $\beta_{e+1}$. Under the beacon-unbiasability assumption, the
event \(\mathcal{Q}_{e+1} \cap \mathcal{C}_{e+1}\) is a uniform random
sample from the size-$(t_{\mathrm{old}})$ subsets of $\mathcal{O}$;
the probability that $\mathcal{A}$'s
$|\mathcal{C}_{e+1}| \le t-1$ corruption covers $\mathcal{Q}_{e+1}$ entirely
is at most
$\binom{t-1}{t}\binom{n}{t}^{-1} = 0$ (corruption budget is strictly less than $t$).
\item By \Cref{thm:reshare-pkinv}, $\mathit{pk}$ is invariant.
\item By the resharing privacy theorem
(\Cref{thm:reshare-priv} below), the joint view of $\mathcal{A}$
across $(\mathcal{C}_e, \mathcal{C}_{e+1})$ is independent of the
master secret $\mathbf{s}_1$.
\item Forgery in epoch $e+1$ reduces to TS-UF in epoch $e+1$ against
the (independently distributed) new committee shares.
\end{enumerate}
The union bound over $E$ epochs gives the stated advantage. \qed
\end{proof}

% =====================================================================
\subsection{Games DKG-SOUND and DKG-PRIV}
\label{sec:game-dkg}
% =====================================================================

\begin{game}[DKG-SOUND]
\label{game:dkg-sound}
$\mathcal{A}$ corrupts up to $t-1$ parties at DKG start. The
challenger runs $\DKG$ honestly across the remaining parties;
$\mathcal{A}$'s corrupted parties act under its control.
$\mathcal{A}$ wins if the output $\mathit{pk}$ is \emph{not} a valid
FIPS~204 public key (i.e., it does not deserialise correctly, or it
does not satisfy the FIPS~204 public-key range checks).
\end{game}

\begin{theorem}[DKG soundness]
\label{thm:dkg-sound}
For every $\PPT$ adversary,
$\Adv^{\mathsf{DKG\textsf{-}SOUND}}_{\mathrm{Pulsar\text{-}M}}(\mathcal{A}) \le \Adv^{\MSIS}_{q, k, 2\ell, \beta_{\mathrm{bind}}}(\mathcal{B}_{\MSIS}) + q_H / 2^{256}$.
\end{theorem}

\begin{proof}[Sketch]
Honest parties verify Pedersen identities at Round~2. A corrupted
party that contributes a malformed share triggers
\textsc{ComplaintBadDelivery}, removing its contribution from the
aggregation. The remaining honest contributions are
$\chi_\eta$-distributed, so $\mathbf{s}_1 = \sum_{i \in \mathcal{H}}
\mathbf{c}_{i,0}$ has a distribution that is statistically close to
$\chi_\eta^\ell$ scaled by $|\mathcal{H}|$ (which we re-randomise to
$\chi_\eta^\ell$ via the FIPS~204 secret-key sampling discipline).
The Pedersen binding bound prevents the adversary from substituting a
different $\mathbf{c}_{i,0}$ between commit and reveal. The remaining
RO term is the random-oracle collision bound for the commit-digest
binding. \qed
\end{proof}

\begin{game}[DKG-PRIV]
\label{game:dkg-priv}
$\mathcal{A}$ corrupts up to $t-1$ parties and runs $\DKG$ honestly;
the challenger flips a bit $b$. If $b = 0$, $\mathcal{A}$ sees the
real DKG transcript including the corrupted parties' shares. If
$b = 1$, the challenger simulates the corrupted parties' shares from
the public $\mathit{pk}$ using a $\Sigma$-protocol-style simulator
that re-randomises the commits without knowledge of the honest
secrets. $\mathcal{A}$ outputs a guess $\hat b$.
\end{game}

\begin{theorem}[DKG share privacy]
\label{thm:dkg-priv}
$\Adv^{\mathsf{DKG\textsf{-}PRIV}}_{\mathrm{Pulsar\text{-}M}}(\mathcal{A}) \le \Adv^{\MLWE}_{q, k, \ell, \eta}(\mathcal{B})$.
\end{theorem}

\begin{proof}[Sketch]
The Pedersen commitments are computationally hiding
(\Cref{thm:pedersen-hiding}). The simulator chooses corrupted-party
shares uniformly from the simulation distribution induced by
$\mathit{pk}$, then programs the commits via the Pedersen-hiding
simulator. Distinguishing real from simulated commits reduces to
$\MLWE$. \qed
\end{proof}

% =====================================================================
\subsection{Games RESHARE-SOUND and RESHARE-PRIV}
\label{sec:game-reshare}
% =====================================================================

\begin{game}[RESHARE-SOUND]
\label{game:reshare-sound}
$\mathcal{A}$ has corrupted up to $t-1$ parties of the old
committee and up to $t-1$ of the new committee. The challenger runs
$\Reshare$ with an honest beacon. $\mathcal{A}$ wins if the new
committee's resulting shares do \emph{not} reconstruct to
$\mathbf{s}_1$ via Lagrange interpolation at $X = 0$.
\end{game}

\begin{theorem}[Reshare soundness]
\label{thm:reshare-sound}
$\Adv^{\mathsf{RESHARE\textsf{-}SOUND}}_{\mathrm{Pulsar\text{-}M}}(\mathcal{A}) \le \Adv^{\MSIS}_{q, k, 2\ell, \beta_{\mathrm{bind}}}(\mathcal{B})$.
\end{theorem}

\begin{proof}[Sketch]
The reshare polynomials $g_i, h_i$ commit to constant terms
$\lambda_i^{\mathcal{Q}_e} \mathbf{s}_i, \lambda_i^{\mathcal{Q}_e} \mathbf{u}_i$
under Pedersen commit. Pedersen binding prevents a deviating
$i \in \mathcal{Q}_e \cap \mathcal{C}$ from substituting a different
constant term. The remaining honest parties' constant terms sum
under Lagrange interpolation to $\mathbf{s}_1$ exactly
(\Cref{thm:reshare-pkinv}). Therefore the master secret is preserved
modulo binding. \qed
\end{proof}

\begin{game}[RESHARE-PRIV]
\label{game:reshare-priv}
$\mathcal{A}$ corrupts up to $t-1$ parties of $\mathcal{O}$ and up to
$t-1$ of $\mathcal{N}$ (allowing distinct corruption sets per epoch,
matching the mobile adversary model). The challenger runs $\Reshare$
honestly with an honest beacon. $\mathcal{A}$ outputs a guess for
$\mathbf{s}_1$. Advantage = probability of correct guess.
\end{game}

\begin{theorem}[Reshare share privacy]
\label{thm:reshare-priv}
For an unbiased beacon $\beta_e$,
$\Adv^{\mathsf{RESHARE\textsf{-}PRIV}}_{\mathrm{Pulsar\text{-}M}}(\mathcal{A}) \le \Adv^{\MLWE}_{q, k, \ell, \eta}(\mathcal{B}) + \Adv^{\mathsf{beacon\textsf{-}unbias}}_{\mathrm{chain}}(\mathcal{B}_{\mathrm{beacon}}) + q^{-1}$.
\end{theorem}

\begin{proof}[Sketch]
The old committee's leaked shares $\{\mathbf{s}_i\}_{i \in
\mathcal{C}_{\mathrm{old}}}$ alone leak no information about
$\mathbf{s}_1$ by Shamir security. The new shares
$\{\mathbf{s}'_j\}_{j \in \mathcal{C}_{\mathrm{new}}}$ similarly leak
nothing on their own. The cross-leakage from the dealing step
contributes $\{g_i(j) : i \in \mathcal{Q}_e, j \in
\mathcal{C}_{\mathrm{new}}\}$. Conditioned on $\mathbf{s}_i$ (which
$\mathcal{A}$ already knows if $i \in \mathcal{C}_{\mathrm{old}}$,
otherwise hidden), the distribution of $g_i(j)$ for $|\{j\}| < t_{\mathrm{new}}$ is uniform in $\Rq^\ell$ by the high-degree
random coefficients of $g_i$; therefore the dealing-step view admits
a simulator that draws the leaked $g_i(j)$ values uniformly without
knowing $\mathbf{s}_1$.

The beacon-unbiasability assumption ensures
$\mathcal{Q}_e \cap \mathcal{C}_{\mathrm{old}}$ is a fresh uniform
sample with respect to $\mathbf{s}_1$, preventing adversarial gaming
of the quorum.

The $q^{-1}$ term is the negligible probability that a Lagrange
denominator collides in the beacon-induced ordering. \qed
\end{proof}

% =====================================================================
\subsection{Game RESTART-NL (rejection-restart non-leakage)}
\label{sec:game-restart-nl}
% =====================================================================

We re-state \Cref{thm:restart-non-leak} from the main text as an
explicit game for completeness and so the reduction to $\KMAC$'s
PRF security is recorded alongside the other games.

\begin{game}[RESTART-NL]
\label{game:restart-nl}
$\mathcal{A}$ corrupts up to $t-1$ parties of a fixed quorum $T$.
The challenger runs $\Sign$ on a fixed honest target $(\mathit{pk}, \mu, T)$.
$\mathcal{A}$ requests $K$ distinct rejection-restart attempts
(distinct $\kappa = 1, \dots, K$); the challenger executes each
attempt and reveals the full per-attempt transcript. $\mathcal{A}$
outputs a guess for $\mathbf{s}_i$ for some honest party $i \in T
\setminus \mathcal{C}$.
\end{game}

\begin{theorem}[Restart non-leakage]
\label{thm:restart-nl-game}
For every $\PPT$ adversary against \Cref{game:restart-nl},
\[
\Adv^{\mathsf{RESTART\textsf{-}NL}}_{\mathrm{Pulsar\text{-}M}}(\mathcal{A})
\;\le\;
K \cdot \Adv^{\mathsf{PRF}}_{\KMAC}(\mathcal{B})
+ \Adv^{\mathsf{TS\textsf{-}UF}}_{\mathrm{Pulsar\text{-}M}}(\mathcal{A}').
\]
\end{theorem}

The proof is the hybrid argument of the main-text
\Cref{thm:restart-non-leak} restated formally. The cross-message
extension --- replacing the fixed $\mu$ with adversarially-chosen
$\mu^{(\kappa)}$ between restarts --- is recorded as open in
\S\ref{sec:knownlimits}.

% =====================================================================
\subsection{Summary of reductions}
\label{sec:sec-summary}
% =====================================================================

\begin{center}
\begin{tabular}{l l}
\toprule
game & reduction target \\
\midrule
TS-UF (static unforgeability) & $\MLWE_{q,k,\ell,\eta}$ + $\MSIS_{q,k,2\ell,\beta_{\mathrm{bind}}}$ + EUF-CMA(ML-DSA) \\
OUT-IND (FIPS~204 interchangeability) & $\MLWE_{q,k,\ell,\eta}$ \\
ROB (identifiable abort) & Pedersen binding ($\MSIS$) + Ed25519 signature soundness \\
ADAPT (adaptive corruption) & open --- deferred to round-2 \\
MOB (mobile-adversary) & TS-UF + RESHARE-PRIV + beacon unbiasability \\
DKG-SOUND & $\MSIS$ + RO collision \\
DKG-PRIV & $\MLWE$ \\
RESHARE-SOUND & $\MSIS$ \\
RESHARE-PRIV & $\MLWE$ + beacon unbiasability \\
RESTART-NL (rejection-restart) & PRF($\KMAC$) + TS-UF \\
\bottomrule
\end{tabular}
\end{center}

The dominant terms across all games are $\MLWE_{q, k, \ell, \eta}$
and ML-DSA EUF-CMA, both of which take the FIPS~204
parameter-pinned values (see \S\ref{sec:parameters}). Pulsar
inherits ML-DSA's security level modulo the threshold-protocol
overhead, which is dominated by the Pedersen-binding $\MSIS$ term
and the random-oracle programming penalty.

\paragraph{Forward-security note.}
The MOB game (\Cref{game:mob}) implies forward security at the epoch
granularity: revealing all shares at epoch $e$ leaks no information
about $\mathbf{s}_1$ in epoch $e' > e$ provided fewer than $t-1$
parties have leaked their epoch-$e'$ shares. Combined with $\Reshare$
at the epoch boundary, this gives the standard
proactive-secret-sharing forward-security guarantee.


% =====================================================================
\section{Concrete parameters}
\label{sec:parameters}
% =====================================================================

% Pulsar concrete parameter sets — included by pulsar.tex §"Concrete parameters"

\subsection*{Parameter sets}

% Three parameter sets, mirroring FIPS 204 ML-DSA's three security levels.
% Numeric values to be filled in from FIPS 204 §4 once the M-LWE port is
% finalized; the Pulsar parameter table is structurally identical to
% ML-DSA's because the underlying ring + module shape are inherited.

\begin{center}
\begin{tabular}{l r r r}
\hline
parameter & Pulsar-44 & Pulsar-65 & Pulsar-87 \\
\hline
NIST PQ category               & TBD (Cat 2) & TBD (Cat 3) & TBD (Cat 5) \\
classical security (bits)      & TBD ($\geq128$) & TBD ($\geq192$) & TBD ($\geq256$) \\
quantum security (bits)        & TBD             & TBD             & TBD             \\
modulus $q$                    & 8380417    & 8380417    & 8380417 \\
ring degree $n$                & 256        & 256        & 256 \\
module dimension $k$           & 4          & 6          & 8 \\
secret dimension $\ell$        & 4          & 5          & 7 \\
secret coefficient bound $\eta$& 2          & 4          & 2 \\
challenge weight $\tau$        & 39         & 49         & 60 \\
hint bound $\omega$            & 80         & 55         & 75 \\
threshold $(n, t)$             & runtime    & runtime    & runtime \\
\hline
signature size (bytes)         & 2420       & 3309       & 4627    \\
public key size (bytes)        & 1312       & 1952       & 2592    \\
\hline
\end{tabular}
\end{center}

% TODO. Fill in security-level numbers from a lattice-estimator run.
% Pin the script in estimator/params.sage and the result table in
% estimator/results.md. Cite the estimator commit hash.

\subsection*{Lattice-estimator inputs}

% TODO. Each parameter set's M-LWE instance, run through the
% lattice-estimator (https://github.com/malb/lattice-estimator) for
% classical and quantum cost estimates. Reproducible from
% estimator/params.sage.

\subsection*{NIST MPTC compliance}

\textbf{Required parameterization} (NIST IR 8214C §4.5): at least one
parameter set with classical $\geq 128$ bits, NIST PQ category $\geq 1$,
statistical $\geq 40$ bits. Pulsar-44 satisfies this.

\textbf{Suggested parameterization}: at least one parameter set with
classical $\geq 192$ bits, NIST PQ category $\geq 3$, statistical
$\geq 64$ bits. Pulsar-65 satisfies this.

Pulsar-87 covers the highest NIST PQ category for completeness.


% =====================================================================
\section{Encodings}
\label{sec:encodings}
% =====================================================================

This section specifies the \emph{structure} of every wire-level
message. Byte-level layouts (offsets, lengths, endianness) freeze at
the encoding-freeze gate (DD-008, end of August 2026). Until that
gate, the structural commitments below are normative; byte offsets
within them are not.

\paragraph{Common conventions.}
\begin{itemize}[leftmargin=*, itemsep=0pt]
\item Every wire message starts with a fixed-length version tag:
$(\text{`PULSAR-1'}, \text{kind\_byte}, \text{epoch\_u32})$.
Tag length: TBD at freeze.
\item Lengths are big-endian unsigned. Length-prefixed for every
variable-size field.
\item Polynomials in $\Rq$ are encoded coefficient-by-coefficient
using the FIPS~204 byte-packing for the per-poly bit width (see
FIPS~204~\S7.1, \texttt{bitlen}). Pulsar re-uses the FIPS~204
packings verbatim for $\mathbf{s}_1, \mathbf{s}_2, \mathbf{z}, \mathbf{w}$.
\item Module vectors are encoded as their entries concatenated in
order. Module matrices are encoded row-major.
\item All multi-party transcripts include the
$(\mathit{pk}, \text{epoch}, \mathit{sid}, T)$ tuple as a
prefix before the message body, bound under a $\cSHAKE$ transcript
hash with the appropriate customisation tag.
\end{itemize}

\paragraph{$\DKG$ Round-1 message.}
\(\{\mathbf{C}_{i,k}\}_{k=0}^{t-1}\) followed by the per-recipient
private channel envelope containing $(f_i(j), g_i(j))$ for each
$j \ne i$. The broadcast portion (commits) is length-pinned because
$t$ and the module shape are known at the start of the protocol.

\paragraph{$\DKG$ Round-1.5 broadcast.}
$(i, H_i)$ where $H_i$ is the 32-byte commit digest.

\paragraph{$\DKG$ Round-2 message.}
$(i, \mathbf{s}_i^{\mathrm{aggregated}}, \mathbf{u}_i^{\mathrm{aggregated}})$
\emph{never appears on the wire}: aggregation is local. The Round-2
artifact is a per-party acceptance signature on the aggregated
$(\mathbf{t}, \mathit{pk})$ value, signed under the party's long-term
identity key.

\paragraph{$\Sign$ Round-1 broadcast.}
$(\mathit{sid}, \kappa, i, D_i, \{\mathrm{MAC}_{i \to j}\}_{j \in T \setminus \{i\}})$.

\paragraph{$\Sign$ Round-2 broadcast.}
$(\mathit{sid}, \kappa, i, \bar{\mathbf{w}}_i, \mathbf{z}_i, \mathbf{r}_i, \text{Round2integrity})$.

\paragraph{$\Sign$ output (final signature).}
The FIPS~204 ML-DSA signature triple $(\tilde c, \mathbf{z}, \mathbf{h})$,
encoded \emph{exactly} per FIPS~204~\S7.2. \textbf{No Pulsar envelope.}

\paragraph{$\Reshare$ messages.}
Structurally identical to $\DKG$ except the commit constant terms
are $\lambda_i^{\mathcal{Q}_e} \mathbf{s}_i$ rather than fresh; and the
broadcast carries the beacon $\beta_e$ to bind the quorum choice.

\paragraph{Complaint records.}
$(\text{kind}\,\text{u8},\,\text{accuser}\,\text{u32},\,\text{accused}\,\text{u32},\,\text{epoch}\,\text{u32},\,\text{evidence-blob}\,\text{var},\,\text{signature}\,\text{64B})$.
Length-pinned per kind.

\paragraph{Activation certificate.}
A FIPS~204 ML-DSA signature on the reshare transcript hash under the
unchanged $\mathit{pk}$. The hash is bound under the
\texttt{PULSAR-RESHARE-TRANSCRIPT-V1} customisation tag.

\paragraph{Open at freeze.}
Outstanding decisions for the encoding freeze:
(i) the exact version tag layout for cross-implementation feature
gating,
(ii) the exact embedding of $\bar{\mathbf{w}}_i$ in the Round-2
broadcast (compressed vs.\ uncompressed FIPS~204
\texttt{w1Encode}),
(iii) the byte layout of the per-party MAC bundle (sparse vs.\
dense indexing),
(iv) the exact ordering and field widths for complaint evidence blobs.

% =====================================================================
\section{System model}
\label{sec:system-model}
% =====================================================================

% Pulsar system model — included by pulsar.tex §"System model"

\section*{System model}

% TODO. Cover:
%
% Network model
%   - synchronous vs asynchronous (Pulsar assumes synchronous broadcast
%     within the round window; partially-synchronous is future work)
%   - broadcast channel: authenticated, eventual delivery
%   - point-to-point channel: authenticated, encrypted
%   - public bulletin: authenticated, append-only (consensus-anchored)
%
% Setup assumptions
%   - PKI: each party has a long-term verification key advertised on the
%     public bulletin (e.g. lux/consensus identity record per HIP-0077)
%   - Random beacon: shared randomness oracle for reshare quorum selection
%     (lux consensus randomness layer)
%
% Trust model
%   - Up to t-1 corrupted parties (static), t-1 per epoch (mobile)
%   - Honest majority of (n - t + 1) needed for liveness
%
% Abort model
%   - Identifiable abort: protocol surfaces a signed complaint when a
%     deviation is detected, with the deviating party ID
%   - Slashing: governance layer consumes complaints (out of scope here)
%
% Preprocessing
%   - DKG runs once per validator-set epoch
%   - Reshare runs at epoch boundary (typical cadence: hours to days)
%   - Sign runs per-message; preprocessing-free if combined with non-cached
%     sampler, or with optional Round-0 mask precomputation per Pulsar's
%     offline/online split
%
% Rotation
%   - Reshare protocol updates shares without changing pk (per HJKY97)
%   - Members may join / leave between epochs subject to (n', t') constraint
%
% Monitoring
%   - Public artifacts: pk, commitments, signed complaints, signature outputs
%   - Private artifacts: shares, blinds, masks, samples — never logged
%
% Deployment surface
%   - Reference assumes pure software
%   - HSM integration is future work (out of scope for round 1 MPTC)


% =====================================================================
\section{Implementation considerations}
\label{sec:impl}
% =====================================================================

\subsection{Constant-time signing}
\label{sec:impl-ct}

All secret-dependent comparisons MUST be constant-time. In
particular:

\begin{itemize}[leftmargin=*, itemsep=0pt]
\item Pedersen-identity verification (\S\ref{sec:dkg} Round~2 and
\S\ref{sec:reshare}) compares full module vectors in $\Rq^k$. The
comparison MUST iterate over every coefficient slot of every entry
unconditionally and reduce the per-slot equality bit into a single
accumulator (e.g.\ \texttt{subtle.ConstantTimeCompare} over a
serialised view, as in
\texttt{pulsar/dkg2/dkg2.go: constTimePolyEqual}).
\item FIPS~204 norm checks ($\|\mathbf{z}\|_\infty < \gamma_1 - \beta$
etc.) MUST be performed with branchless arithmetic. The
\texttt{pulsar/sign/sign.go: CheckL2Norm} implementation is the
reference (modulo Pulsar's switch from L2 norm to L-infty per
FIPS~204).
\item Rejection-restart MUST NOT short-circuit on the first failing
norm: all norms are evaluated, then the accept/restart decision is
made from the accumulated result. This prevents timing leakage of
\emph{which} bound failed.
\item No \texttt{log.*}, \texttt{fmt.Print*}, or panic in any
secret-touching function (DD-007). The Pulsar reference's
pre-fix \texttt{log.Println(Bsquare, sumSquares)} (HIP-0077 F8) is
exactly the failure mode this rule prevents; the Pulsar reference
implementation will be CI-gated against any logging call inside
\texttt{ref/go/sign/}, \texttt{ref/go/dkg/}, \texttt{ref/go/reshare/}.
\end{itemize}

\subsection{Rejection-sampling restart without secret leakage}
\label{sec:impl-rejection-restart}

This is the hardest cryptographic engineering problem in Pulsar.
FIPS~204 ML-DSA's rejection-with-restart structure (\S\ref{sec:prelim-rejection})
interacts with threshold signing in three load-bearing places.

\paragraph{The threat.}
A na\"ive port of Pulsar's Round-1 PRNG keying would seed party
$i$'s PRNG from the secret share $\mathbf{s}_i$ alone:
\(
\text{seed}_i \gets \text{Hash}(\mathbf{s}_i)
\).
Two distinct sign attempts with the same $(\mathit{pk}, \mu)$ would
then sample byte-identical $\mathbf{y}_i$ values, and an attacker
who saw the resulting $\mathbf{z}$ values from two attempts could
compute
\[
\mathbf{z}^{(1)} - \mathbf{z}^{(2)} = c^{(1)} \mathbf{s}_1 - c^{(2)} \mathbf{s}_1
= (c^{(1)} - c^{(2)}) \mathbf{s}_1,
\]
and \emph{recover the secret} by dividing by the (low-norm,
small-coefficient) difference $c^{(1)} - c^{(2)}$ over $\Rq$. This is
the CRIT-1 vulnerability flagged in Pulsar's pre-fix code
(\texttt{pulsar/primitives/hash.go: PRNGKey}) and addressed
2026-05-03 in coordination with the C++ port at
\texttt{luxcpp/crypto}. The fix is the
\texttt{PRNGKeyForRound} function: the PRNG seed mixes
$(\mathbf{s}_i, \mathit{sid})$ so distinct sessions get distinct
seeds.

\paragraph{The Pulsar domain separation.}
Pulsar extends the Pulsar fix to handle three orthogonal
domains:

\begin{enumerate}[leftmargin=*, itemsep=0pt]
\item \textbf{Across signatures} (distinct $\mathit{sid}$). Each
sign attempt gets a fresh $\mathit{sid}$, monotonically advanced
by the protocol layer above (Lux consensus uses the block height
plus the within-block signature index as $\mathit{sid}$).
\item \textbf{Across rejection-restart attempts within one
$\mathit{sid}$} (distinct $\kappa$). Each restart advances the
attempt counter. The per-attempt PRNG seed is derived as:
\[
\text{seed}_i^{(\kappa)} \;=\;
\KMAC(\,K_i^{\mathrm{prng}},\,
(\mathit{sid}, \kappa, T, i, \mathit{pk}, \mu),\,
256,\,
\text{`PULSAR-SIGN-PRNG-V1'}\,),
\]
where $K_i^{\mathrm{prng}}$ is a per-party PRNG key derived at DKG
time from $\mathbf{s}_i$ via $\cSHAKE$ with the customisation tag
\texttt{`PULSAR-SIGN-PRNGKEY-V1'}. Both $\mathit{sid}$ and $\kappa$
are bound into the seed.
\item \textbf{Across quorums on the same $\mathit{sid}$} (distinct
$T$). If two distinct quorums coordinate on the same $\mathit{sid}$
(e.g.\ a network split followed by reconciliation), the seed must
diverge so that the two attempts use independent masks. The
quorum set $T$ is bound into the seed as above.
\end{enumerate}

\paragraph{Cross-restart state.}
Across a rejection-sampling restart boundary, only the following
state is preserved:
\begin{itemize}[leftmargin=*, itemsep=0pt]
\item the long-term shares $(\mathbf{s}_i, \mathbf{u}_i, K_i^{\mathrm{prng}}, \{K_{i,j}\})$,
\item the public coordination state $(\mathit{sid}, \mathit{pk}, T, \mu)$,
\item the attempt counter $\kappa$.
\end{itemize}
\textbf{Everything else} is recomputed: PRNG seeds, masks
$\mathbf{y}_i$, commits $\mathbf{w}_i$, MACs, per-party response
contributions $(\mathbf{z}_i, \mathbf{r}_i)$. There is no
``cached'' per-round value that crosses the boundary. This
discipline is the cryptographic justification for the next theorem.

\begin{theorem}[Restart non-leakage]
\label{thm:restart-non-leak}
Let $\mathcal{A}$ be an adversary that observes a transcript of $K$
rejection-sampling attempts on the same $(\mathit{pk}, \mu, T)$
(distinct $\kappa = 1, \dots, K$) and outputs a guess
$\hat{\mathbf{s}}_i$ for some honest party $i \in T \setminus
\mathrm{corrupt}$. Under the assumption that $\KMAC$ is a PRF in its
key, $\mathcal{A}$'s advantage in recovering $\mathbf{s}_i$ from $K$
attempts is at most $K \cdot \Adv^{\mathrm{prf}}_{\KMAC}(\mathcal{B})$
plus the single-attempt advantage:
\[
\Pr[\mathcal{A}(\mathrm{transcript}_{1..K}) = \mathbf{s}_i]
\;\le\;
K \cdot \Adv^{\mathrm{prf}}_{\KMAC}(\mathcal{B})
+ \Pr[\mathcal{A}'(\mathrm{transcript}_{1}) = \mathbf{s}_i],
\]
for some efficient $\mathcal{B}$ and an equally-capable
$\mathcal{A}'$ restricted to a single transcript.
\end{theorem}

\begin{proof}[Proof sketch]
The PRF assumption on $\KMAC$ implies that
$\{\text{seed}_i^{(\kappa)}\}_{\kappa=1}^K$ is computationally
indistinguishable from $K$ independent uniform samples in
$\{0,1\}^{256}$, with distinguishing advantage at most
$\Adv^{\mathrm{prf}}_{\KMAC}(\mathcal{B})$ per call (i.e.\
total $\le K \cdot \Adv^{\mathrm{prf}}_{\KMAC}(\mathcal{B})$ by a
standard hybrid). Once we condition on $K$ independent seeds, the
$K$ masks $\{\mathbf{y}_i^{(\kappa)}\}_{\kappa=1}^K$ are independent
uniform-interval samples, and the $K$ responses
$\{\mathbf{z}_i^{(\kappa)}\}$ are $K$ independent functions of
$\mathbf{s}_i$ each with the same distribution as a single
honest-party response. The adversary thus gains no more information
than from $K$ \emph{independently sampled} signatures, which is the
same as from a single signature plus $K-1$ ``free'' instances under
M-LWE. The single-attempt advantage bound completes the
inequality. \qed
\end{proof}

\begin{remark}[Stronger statement deferred]
A stronger statement would replace the right-hand side
single-attempt advantage with the EUF-CMA advantage of FIPS~204
against $K$-fold signing queries on the same message. This is
straightforward in the single-message setting but requires a careful
hybrid for cross-message correlations (e.g.\ if $\mathcal{A}$
chooses adaptive messages between restarts). We mark this as
Round-2-MPTC work in \S\ref{sec:knownlimits}.
\end{remark}

\paragraph{Implementation discipline.}
The Pulsar reference implementation enforces the cross-restart
state discipline structurally:

\begin{itemize}[leftmargin=*, itemsep=0pt]
\item The per-attempt state is a fresh struct allocated at the top
of each restart iteration. No cross-iteration field carries forward.
\item The PRNG seeding routine takes
$(\mathit{sid}, \kappa, T, i, \mathit{pk}, \mu)$ as inputs and is
the unique source of per-attempt randomness. Any attempt to use
\texttt{crypto/rand} inside the inner loop is a CI failure.
\item Constant-time discipline (\S\ref{sec:impl-ct}) is preserved
across restarts: the restart decision itself is constant-time in
the failure mode (which norm check failed is leaked only via the
fact of failure, not via timing).
\end{itemize}

\subsection{Performance characteristics}
\label{sec:impl-perf}

Concrete benchmarks defer to the experimental-evaluation report
(\texttt{bench/results/REPORT.md}, frozen at submission). Pulsar's
expected complexity scales as follows for parameter set ML-DSA-65,
$(n, t) = (7, 4)$:

\begin{center}
\begin{tabular}{l l l l}
\toprule
operation & messages & local compute & per-party state \\
\midrule
$\DKG$ (3 rounds)
& $O(n^2)$ broadcast + $O(n)$ per recipient
& $O(n \cdot k \ell)$ per party
& $O(t \cdot k \ell)$ \\
$\Sign$ (2 rounds, no restart)
& $O(t^2)$ MACs + $O(t)$ broadcast
& $O(k \ell)$ per party
& $O(\ell)$ \\
$\Sign$ (avg.\ FIPS~204 restart count $\approx 5$)
& $5 \times$ above
& $5 \times$ above
& unchanged \\
$\Reshare$ (3 rounds)
& $O(t_{\mathrm{old}} \cdot |\mathcal{N}|)$ private + $O(n)$ broadcast
& $O(t_{\mathrm{old}} \cdot k \ell)$
& $O(t_{\mathrm{new}} \cdot k \ell)$ \\
\bottomrule
\end{tabular}
\end{center}

\paragraph{Signature size.} The output is a single FIPS~204
signature: ML-DSA-44 (2420 B), ML-DSA-65 (3309 B), ML-DSA-87
(4627 B). \emph{No threshold envelope.}

\subsection{Side-channel resistance}
\label{sec:impl-sidechannel}

The reference implementation targets constant-time semantics as
specified in \S\ref{sec:impl-ct}. Physical side-channel resistance
(power analysis, EM, fault injection, cache timing) is out of scope
for the round-1 MPTC submission; see
\texttt{docs/known-limitations.md} for the deferred mitigations and
\S\ref{sec:knownlimits} for the round-2 roadmap.

% =====================================================================
\section{Test vectors}
\label{sec:vectors}
% =====================================================================

The KAT suite (frozen at submission) lives in
\texttt{vectors/kat-v1.\{json,rsp\}}. Each entry encodes a complete
ceremony:

\begin{itemize}[leftmargin=*, itemsep=0pt]
\item Setup: $(q, n_{\mathrm{ring}}, k, \ell, \eta, \tau, \gamma_1, \gamma_2, \omega, \beta, d, \lambda, n, t)$.
\item DKG: per-party seeds (deterministic), Round-1 outputs (commits and dealt shares), Round-2 verification results, final $\mathit{pk}$.
\item Sign: $(\mathit{sid}, T, \mu, K_i^{\mathrm{prng}})$, per-party Round-1 and Round-2 transcripts, per-attempt $\kappa$ trace, final $\sigma$, FIPS~204 \texttt{ML-DSA.Verify} result.
\item Reshare: old shares, new committee, beacon $\beta_e$, per-party transcripts, new shares, activation certificate, $\mathit{pk}$-invariance proof.
\end{itemize}

Cross-implementation KAT replay: the Go reference (\texttt{ref/go/})
and the eventual C reference (\texttt{ref/c/}) MUST produce
byte-equal outputs from the same seeded inputs. The KAT generation
oracle is locked at \texttt{cmd/pulsar-oracle}.

% =====================================================================
\section{Complexity analysis}
\label{sec:complexity}
% =====================================================================

\paragraph{Communication.}
$\DKG$: $O(n^2 \cdot t \cdot k)$ bits broadcast (commits) plus
$O(n^2 \cdot t \cdot \ell)$ bits private (shares and blinds).
$\Sign$: $O(t^2)$ MAC bits plus $O(t \cdot (k + \ell))$ bits
broadcast per round. $\Reshare$: $O(t_{\mathrm{old}} \cdot
(t_{\mathrm{new}} \cdot k + |\mathcal{N}| \cdot \ell))$.

\paragraph{Computation.}
$\DKG$: $O(n \cdot t \cdot k\ell)$ ring operations per party.
$\Sign$: $O(k\ell)$ ring operations per party per round (one matrix
multiplication, plus the Lagrange-scalar response). The aggregator
performs an additional $O(t \cdot \ell)$ sum.
$\Reshare$: $O(t_{\mathrm{old}} \cdot t_{\mathrm{new}})$ per old quorum
member, $O(t_{\mathrm{old}})$ per new committee member.

\paragraph{Memory.}
Per-party state: $O(t \cdot k\ell)$ (shares plus DKG commitment
backlog). $\Sign$ holds an $O(\ell)$ working set within the inner
loop.

\paragraph{Round count.}
$\DKG$: 3. $\Sign$: 2 per attempt; the FIPS~204 acceptance
probability is approximately $1 - e^{-\beta\tau/2\gamma_1}$ per
attempt (FIPS~204~\S6.1.1), giving an expected $\sim 2.4$ attempts at
ML-DSA-65. $\Reshare$: 3.

% =====================================================================
\section{Deployment considerations}
\label{sec:deployment}
% =====================================================================

\paragraph{Network model.}
The protocols assume synchronous broadcast within the round window:
all honest parties receive all honest broadcast messages within a
bounded $\Delta$. Partially synchronous (eventual delivery within a
sufficient round window) is acceptable in practice; full
asynchrony is not within scope for round-1 MPTC.

\paragraph{Abort and recovery.}
On any complaint that reaches the disqualification threshold, the
chain layer removes the disqualified party from the validator set
and triggers a $\Reshare$ to a new committee. Liveness is preserved
provided fewer than $t-1$ parties are disqualified per epoch.

\paragraph{Rotation cadence.}
Production deployments are expected to run $\Refresh$ every
$\sim$ hour (background hygiene against mobile adversary) and
$\Reshare$ at every validator-set change (typically hours to days
between epochs).

\paragraph{Monitoring.}
Public artifacts safe for observability: $\mathit{pk}$, commitments,
signed complaints, signature outputs, beacon values, per-party
identity-key signatures. Private artifacts that MUST NOT be logged:
shares, blinds, masks, samples, PRNG seeds, per-round derived state,
norm-check intermediate values.

\paragraph{HSM integration.}
The reference implementation is software. HSM integration --- where
each party's share lives in an HSM and the HSM exposes only the
$\Sign$ Round-2 contribution --- is out of scope for round-1 MPTC
and is the canonical roadmap item for round-2 hardening.

% =====================================================================
\section{Comparison with related work}
\label{sec:related}
% =====================================================================

\paragraph{Pulsar (R-LWE), Class~S~\cite{boschini2024corona}.}
The companion Lux scheme over $\Rq$, single polynomial. Faster,
simpler, but \emph{not} output-interchangeable with any NIST primitive.
Different wire format. Pulsar shares Pulsar's protocol structure
(2-round signing, Pedersen DKG, HJKY97 reshare) lifted from $\Rq$ to
$\Rq^k$.

\paragraph{GKS24~\cite{gks24}, PQCrypto 2024.}
Two-round threshold lattice signatures from threshold homomorphic
encryption (G\"ur--Katz--Silde). Output is a TFHE-shaped artefact, not
ML-DSA. Class~S. Higher signature size than ML-DSA but tighter
proofs.

\paragraph{DOTT21~\cite{dott21}, PKC 2021.}
The foundational Damg\aa rd--Orlandi--Takahashi--Tibouchi two-round
$n$-out-of-$n$ multisignature from lattices. Inspires the
commit-and-aggregate structure of Pulsar/Pulsar. Restricted to
$n$-out-of-$n$; Pulsar extends to $t$-out-of-$n$ via Shamir.

\paragraph{Olingo~\cite{olingo2025}, ePrint 2025/1789.}
Concurrent threshold ML-DSA candidate. Comparison points: round
count, output interchangeability claim, DKG soundness, side-channel
posture. Full comparison table deferred to the experimental
evaluation report; structural comparison: Olingo and Pulsar both
target Class N1, both use $\Rq^k$, with differing rejection-restart
strategies.

\paragraph{THED~\cite{thed2026}, ePrint 2026/638.}
Threshold extension of FIPS~204. Comparison deferred to round-1
preview; THED was not public at Pulsar's draft cut-off.

\paragraph{HJKY97~\cite{hjky97}.}
The foundational proactive-resharing model. Pulsar's $\Reshare$
protocol is the natural lift of the HJKY97 construction to the M-LWE
setting, with the beacon-randomised quorum selection (DD-006) as the
Pulsar-specific strengthening for mobile-adversary attacks.

% =====================================================================
\appendix
\section{Notation summary}
\label{sec:notation-table}
% =====================================================================

\begin{center}
\begin{tabular}{l l}
\toprule
symbol & meaning \\
\midrule
$q$ & prime modulus, $q = 2^{23} - 2^{13} + 1 = 8\,380\,417$ \\
$n_{\mathrm{ring}}$ & ring degree, $n_{\mathrm{ring}} = 256$ \\
$\Rq$ & $\Zq[X]/(X^{n_{\mathrm{ring}}} + 1)$ \\
$(k, \ell)$ & module dimensions, FIPS~204 dependent \\
$(\eta, \tau, \gamma_1, \gamma_2, \omega, \beta, d, \lambda)$ & FIPS~204 parameter set \\
$(n, t)$ & threshold parameters, $n \ge 2t-1$ \\
$\mathbf{A} \in \Rq^{k \times \ell}$ & public matrix, derived from $\rho$ via FIPS~204 \texttt{ExpandA} \\
$\mathbf{B} \in \Rq^{k \times \ell}$ & Pedersen blinding matrix, derived via \texttt{PULSAR-EXPANDB-V1} \\
$\mathbf{s}_1 \in \Rq^\ell$ & master secret, $\chi_\eta$-distributed \\
$\mathbf{u} \in \Rq^\ell$ & Pedersen blinding, $\chi_\eta$-distributed \\
$\mathbf{s}_2 \in \Rq^k$ & FIPS~204 noise, $= \mathbf{B}\mathbf{u}$ in Pulsar \\
$\mathbf{t} = \mathbf{A}\mathbf{s}_1 + \mathbf{B}\mathbf{u}$ & FIPS~204 public-key inner vector (Pedersen-shaped) \\
$(\mathbf{t}_1, \mathbf{t}_0) = \Power(\mathbf{t}, d)$ & FIPS~204 public/private decomposition \\
$\mathit{pk} = (\rho, \mathbf{t}_1)$ & FIPS~204 public key \\
$(\mathbf{s}_i, \mathbf{u}_i)$ & party $i$'s Shamir share pair \\
$\lambda_i^T$ & Lagrange coefficient at $X = 0$ for $i \in T$ \\
$\mathbf{y}_i$ & party $i$'s mask, uniform interval $U_{\gamma_1}^\ell$ \\
$\mathbf{w}_i = \mathbf{A}\mathbf{y}_i$ & party $i$'s commit, in $\Rq^k$ \\
$\bar{\mathbf{w}}_i = \HighBits(\mathbf{w}_i, 2\gamma_2)$ & high-bits commit \\
$c \in \Rq$ & FIPS~204 challenge, $\texttt{SampleInBall}(\tilde c)$ \\
$\mathbf{z} \in \Rq^\ell$ & aggregated response, $\mathbf{y} + c\mathbf{s}_1$ \\
$\mathbf{h}$ & FIPS~204 hint vector \\
$\sigma = (\tilde c, \mathbf{z}, \mathbf{h})$ & FIPS~204 signature triple \\
$\mathit{sid}$ & per-attempt session id \\
$\kappa$ & per-attempt rejection-restart counter \\
$T \subseteq \{1, \dots, n\}$ & quorum, $|T| = t$ \\
$\mathcal{Q}_e$ & beacon-selected reshare quorum at epoch $e$ \\
$\beta_e$ & chain randomness beacon at epoch $e$ \\
\bottomrule
\end{tabular}
\end{center}

\section{Patent claims}
\label{sec:patent}
Final patent disclosure assembled from
\texttt{docs/patent-notes-draft.md} at the submission freeze. Lux
Industries Inc.\ asserts no patent claims on the cryptographic body
of Pulsar; the threshold protocol is a port of the publicly-known
Corona~\cite{boschini2024corona} construction with the beacon-quorum
strengthening of DD-006, both of which are explicitly placed into
the public domain by this submission.

\section{Production go-live blockers and deferred theorems}
\label{sec:knownlimits}
Mirror of \texttt{BLOCKERS.md} (replaces the earlier
\texttt{docs/known-limitations.md} framing, which was too lenient).
A 4-red-agent + 1-scientist adversarial audit (2026-05) confirmed
that the v0.1 reference implementation is paper-track-shippable but
\emph{not} production-go-live ready as a strict-PQ threshold signer
in an open public permissionless blockchain. The findings below are
production blockers, not informational caveats; each is filed against
specific files and lines in the \texttt{BLOCKERS.md} document.

\paragraph{Cryptographic gaps that must close before production
deployment.}

\begin{itemize}[leftmargin=*, itemsep=0pt]
\item \textbf{DKG commit-opening (CR-6)} [CLOSED in v0.2.0]: The
unopened \texttt{cSHAKE256(c\_i $\Vert$ blind\_i)} commit was
dropped (Path A in BLOCKERS.md). The v0.1 reference implementation
now realises the Shamir+sum protocol with no separate commit-and-open
layer; binding comes from Round-2 envelope-digest agreement over the
ordered Round-1 broadcast set (\S\ref{sec:dkg}).
\item \textbf{Sign-round session-key derivation (CR-7)} [CLOSED in
v0.2.0]: \texttt{deriveMACKey} was deleted. The threshold signer now
consumes a per-pair ephemeral session key established before Round~1
via authenticated ML-KEM-768 key agreement, where each party
encapsulates a fresh shared secret to the peer's long-term ML-KEM-768
identity public key and signs the ciphertext with their long-term
ML-DSA-65 identity key. The final session key mixes both parties'
contributions through HKDF-SHA3-256 over a canonical-ordered NodeID
pair (\S\ref{sec:sign} and \S\ref{sec:system-model}).
\item \textbf{DKG envelope confidentiality (CR-8)} [CLOSED in v0.2.0]:
Round-1 envelopes are ML-KEM-768-wrapped against each recipient's
long-term identity public key. The wire shape is
$\textsc{DKGShareEnvelope} = (\mathit{KEMCT}, \mathit{Sealed})$ where
$\mathit{KEMCT}$ is the ML-KEM-768 ciphertext (1088 B) and
$\mathit{Sealed}$ is the 128-byte sealed payload carrying
$(\mathit{share} \| \mathit{contribution} \| \mathit{tag})$ XORed with
a cSHAKE256 stream keyed by the KEM-derived envelope key. A passive
network observer learns only ciphertexts; no per-recipient share
leaks (\S\ref{sec:dkg}).
\item \textbf{Adaptive corruption proof (deferred).} The
\Cref{thm:reshare-pkinv} unforgeability theorem is stated against the
static-corruption model only. The adaptive-corruption variant
(Game ADAPT) is the only deferred theorem in the round-1 security
analysis; production deployment under nation-state corruption requires
discharging Game ADAPT.
\item \textbf{Cross-domain isolation between Pulsar (M-LWE) and
Corona (R-LWE).} These are within the same algebraic-lattice
hardness family, not independent failure modes. Defense-in-depth
claims require a hash-based (SLH-DSA) layer also wired into the
QuasarCert envelope.
\item \textbf{Verify constant-time.} \Cref{sec:verify} asserts but
does not measure constant-time behaviour of the verifier path. The
\texttt{ct/dudect/} harness is scaffolded; results are pinned at
the submission tag from a quiet benchmarking machine.
\end{itemize}

\paragraph{Round-1 deferred items (out of scope for v0.1).}

\begin{itemize}[leftmargin=*, itemsep=0pt]
\item \textbf{Asynchronous network model.} Round-1 assumes synchronous
broadcast within the round window. Partial synchrony is acceptable
in practice; full asynchrony requires structural protocol changes
(threshold encryption of round-1 commits is one direction).
\item \textbf{HSM integration.} Software-only reference; HSM-backed
share custody is a round-2 hardening item.
\item \textbf{Physical side-channel resistance.} Constant-time
arithmetic is in scope; physical side channels (power, EM, fault) are
out of scope for round-1.
\item \textbf{Cross-message restart hybrid.} \Cref{thm:restart-non-leak}
proves the restart non-leakage statement for a single $(\mathit{pk},
\mu)$. A stronger cross-message-adaptive hybrid is round-2 work.
\item \textbf{Lattice estimator parameter pins.} Concrete bit-security
estimates from the lattice estimator (\texttt{estimator/}) are
pinned at submission; the underlying estimator commit hash is
recorded in \texttt{estimator/README.md}.
\end{itemize}

\begin{openq}[Adaptive corruption tightness]
The static-to-adaptive corruption gap for two-round threshold
lattice signatures is an active research direction. The
\cite{boschini2024corona} paper's adaptive variant uses a complexity
leveraging argument that loses a multiplicative
$\binom{n}{t}$ factor in the security parameter. Can a tighter
adaptive reduction be constructed for the module-shaped
Pulsar variant without architectural changes? \emph{Unresolved.}
\end{openq}

\begin{openq}[Restart cross-message hybrid]
\Cref{thm:restart-non-leak} proves restart non-leakage for a single
$(\mathit{pk}, \mu)$. A full cross-message-adaptive hybrid would
extend to the case where the adversary chooses $\mu^{(\kappa+1)}$
based on the transcript of attempts on $\mu^{(\kappa)}$. We
conjecture this holds under M-LWE without structural changes, but
the formal hybrid is round-2 work. \emph{Unresolved.}
\end{openq}

\begin{openq}[Beacon-quorum binding for proactive resharing]
The DD-006 beacon-randomised quorum selection assumes the chain
beacon is unbiasable against the adversary. The Lux Quasar
consensus randomness layer provides this property; the formal
binding to the Pulsar reshare MOB game requires a separate
hybrid argument bridging the consensus randomness assumption to
the Pulsar MOB game. \emph{Sketched in security-games.tex; full
proof round-2.}
\end{openq}

\bibliographystyle{plain}
\bibliography{references}

\end{document}
